\section{ELETTROSTATICA DELLE CARICHE FISSE}

\subsection{La legge di Coulomb}

Iniziare oggi, nell'era dell'elettronica più spinta, il corso di elettromagnetismo citando le palline di sambuco e le bacchette di ebanite o di vetro strofinate con pelle di coniglio per convincerci dell'esistenza di fenomeni ``elettrici'', legati ad un particolare stato dei corpi (corpi ``elettrizzati'' o elettricamente carichi), sarebbe -- a dir poco -- anacronistico: la costituzione atomica della materia era già nota alla fine del secolo scorso (fine del 1800), come pure il fatto che gli atomi -- normalmente elettricamente inerti (neutri) -- contengono particelle cariche che possono essere estratte da essi con diverse tecniche. Thomson misurò il rapporto fra la loro carica e la loro massa e trovò che il risultato era indipendente sia dalla tecnica usata per estrarre queste cariche (effetto termoionico, effetto fotoelettrico, raggi X ecc.), sia dalla sostanza da cui venivano estratte: a queste particelle, presenti in tutti gli atomi, fu dato il nome di elettroni. Più tardi (1909) Millikan ne misurò la carica, confermando l'identità di queste particelle. Quindi, quello che gli antichi chiamavano ``fuoco sottile'', qualcosa che conferiva ai corpi quella misteriosa proprietà di attrarre o respingere altri corpi nelle stesse condizioni, altro non è che un eccesso o un difetto di elettroni rispetto allo stato normale del corpo, e che rende il corpo stesso elettricamente attivo, cioè carico.\\
L'esistenza di due tipi di cariche era già nota dai tempi delle bacchette di vetro o di ebanite: per dirla in termini più moderni, era noto dall'esperienza che la proprietà di due cariche di respingersi è una relazione di equivalenza. Fu Franklin, verso la metà del '700 che propose di chiamare positive le cariche di una classe di equivalenza (quella delle cariche delle bacchette di vetro), e negative quelle dell'altra classe: qualunque altro nome sarebbe andato altrettanto bene, e anche la convenzione opposta (e in tal caso gli elettroni risulterebbero cariche positive).\\
Il punto di partenza di tutto l'elettromagnetismo è la legge di interazione fra due cariche, stabilita sperimentalmente da Coulomb nel 1785: la forza che si esercita fra due corpi carichi in quiete, di dimensioni piccole rispetto alla loro distanza, è proporzionale alle due cariche, inversamente proporzionale al quadrato della loro distanza ed ha la direzione della congiungente i due corpi. Se usiamo la convenzione di Franklin ($q_1\cdot q_2>0$ per cariche equivalenti, $q_1\cdot q_2<0$ altrimenti), la legge di Coulomb può essere scritta
\begin{equation}
\vec F_{12} = k\,\frac{q_1q_2}{r_{12}^2}\,\hat r_{12} = k\,\frac{q_1q_2}{|\vec r_1-\vec r_2|^3}\,(\vec r_1-\vec r_2)
\qquad\qquad (\vec r_{12}\equiv\vec r_1-\vec r_2)
\tag{1.1}
\end{equation}
dove $\vec F_{12}$ è la forza che il corpo di carica $q_2$ esercita sul corpo di carica $q_1$, $\hat r_{12}$ è il versore diretto da $q_2$ a $q_1$ e $k$ una costante di proporzionalità positiva.\\
La proporzionalità inversa al quadrato della distanza è fondamentale: l'elettromagnetismo sarebbe completamente diverso da quello che è se la dipendenza da $r$ nella (1.1) fosse diversa (p.es.\ $e^{-\mu r}/r^2$). Tanto per fare un esempio, la velocità delle onde elettromagnetiche (nel vuoto) dipenderebbe dalla loro lunghezza d'onda (il fotone avrebbe massa): in effetti oggi le misure delle differenze dei tempi di arrivo di onde elettromagnetiche di frequenze diverse (luce, radio onde, ecc.)\ provenienti dall'esplosione di stelle lontane costituiscono la prova più convincente della validità della legge $r^{-2}$ alle grandi distanze. Il grado di accuratezza con cui questa legge è stata verificata dipende dalle dimensioni esplorate, in quanto ne dipendono i metodi di verifica: da una parte su $10^9$ per distanze dell'ordine di qualche decimetro a $1/10^6$ per distanze cosmiche.\\
La verifica sperimentale della proporzionalità di $\vec F$ a ciascuna delle due cariche $q_1$ e $q_2$ è sostanzialmente una verifica dell'additività delle cariche (con la convenzione di Franklin di attribuire segni opposti a cariche non equivalenti): la forza che $q_2$ esercita su un (piccolo) corpo costituito da due corpi di cariche $q_1'$ e $q_1''$ (eventualmente uguali) è uguale alla forza che si esercita su un corpo di carica $q_1=q_1'+q_1''$.\\
La costante $k$ dipende dall'unità di misura della carica. A questo proposito si possono seguire due strade:
\begin{itemize}
\item[i)] considerare la carica elettrica, o -- se risulta più comodo per ragioni di riproducibilità -- altra grandezza fisica legata ad essa (p.es.\ la corrente elettrica) come una nuova grandezza fondamentale, fissare un campione che ne definisce l'unità di misura, e quindi determinare (sperimentalmente) il valore di $k$;
\item[ii)] fissare (arbitrariamente) il valore di $k$ nella (1.1), considerata come una costante adimensionata (numero puro); a questo punto la carica elettrica risulta una grandezza derivata le cui dimensioni sono $[Q]=[M]^{1/2}[L]^{3/2}\,[T]^{-1}$, e l'unità di misura è fissata dalla legge di Coulomb.
\end{itemize}
La prima possibilità è quella adottata nell'ambito del SI, in cui viene presa come nuova grandezza fondamentale l'intensità di corrente e se ne definisce l'unità di misura, l'\texttt{ampère} \texttt{A} (del quale daremo la definizione più avanti). Di conseguenza resta fissata l'unità di misura della carica, il \texttt{coulomb} (\texttt{C}), e $k$ (che quindi è una grandezza dimensionata) risulta sperimentalmente $k\approx9\times10^9$. In questo caso si preferisce porre $k\equiv1/4\pi\epsilon_0$, e $\epsilon_0$ è detta \emph{permittività del vuoto}.\\
Quindi nel SI la legge di Coulomb si scrive
\begin{equation}
\vec F_{12}=\frac{1}{4\pi\epsilon_0}\,\frac{q_1q_2}{r_{12}^2}\,\hat r_{12}\,;\qquad
\epsilon_0=8.85\times10^{-12}\ \text{unità del SI (MKSA)}.
\tag{1.1'}
\end{equation}
La carica dell'elettrone (misurata da Millikan) è
\begin{equation}
e=1.6\times10^{-19}\,\mathrm{C}
\tag{1.2}
\end{equation}
(con $e$ si indica il valore assoluto della carica dell'elettrone). Il coulomb corrisponde quindi alla carica di circa $6\times10^{18}$ elettroni: è una carica molto grande, come si capisce dal fatto che fra due cariche di un coulomb alla distanza di 1\,m si esercita una forza di $9\times10^9$\,N, cioè circa $10^9\,\mathrm{kg_p}$!\\
Il sistema di Gauss si avvale invece della seconda possibilità: esso è definito dalle unità del sistema CGS e dalla scelta $k=1$; l'unità di carica elettrica -- il \texttt{\textbf{franklin}} -- è quindi la carica che posta a distanza di 1\,cm da una carica identica la respinge con una forza di 1\,dyna. Il sistema di Gauss nell'ambito dell'elettrostatica è anche detto sistema CGS elettrostatico, e le unità delle grandezze elettriche (alle quali in molti casi non viene dato un nome) vengono sovente indicate con ues (unità elettrostatiche).\\
Quindi nel sistema di Gauss la legge di Coulomb si scrive
\begin{equation}
\vec F_{12}=\frac{q_1q_2}{r_{12}^2}\,\hat r_{12}\,;\qquad\qquad \text{(sistema di Gauss).}
\tag{1.1''}
\end{equation}
Dalla (1.1') risulta che due cariche di 1\,coulomb alla distanza di 1\,cm si respingono con una forza $F=9\times10^9\cdot10^4\,\mathrm{N}=9\times10^{18}\,\mathrm{dyne}$, da cui
\begin{equation}
1\,\mathrm{C}=3\times10^9\,\mathrm{ues}
\tag{1.3}
\end{equation}
e quindi per l'elettrone
\begin{equation}
e=4.8\times10^{-10}\,\mathrm{ues}\,.
\tag{1.2'}
\end{equation}
Nell'Elettromagnetismo (contrariamente alla Meccanica) è necessario specificare il sistema di unità che intendiamo utilizzare in quanto, come abbiamo visto, il SI ed il sistema di Gauss non differiscono semplicemente per la scelta delle unità di misura (come invece l'MKS e il CGS nell'ambito della Meccanica), e nei due sistemi le formule fondamentali vengono scritte in maniera diversa. In questo primo capitolo, tuttavia, in cui l'unica legge fondamentale è la legge di Coulomb e nel quale verranno stabilite le relazioni fondamentali dell'elettrostatica, possiamo rimandare la scelta del sistema di unità utilizzando la (1.1) senza specificare il valore di $k$: tutte le relazioni che dedurremo da essa possono essere immediatamente scritte nel SI o nel sistema di Gauss semplicemente ponendo $k=1/4\pi\epsilon_0$ oppure $k=1$. Quando si introduce il campo magnetico le cose non sono più così semplici in quanto la definizione stessa di campo magnetico è diversa nei due sistemi, come appare dal modo di scrivere la forza che un campo magnetico esercita su una carica in moto: $\vec F=q\vec v\wedge\vec B$ nel SI, e $\vec F=q(\vec v/c)\wedge\vec B$ nel sistema di Gauss ($c$ è la velocità della luce), e di conseguenza tutte le relazioni in cui interviene il campo magnetico differiscono per il fattore $k/c$. Da un punto di vista teorico preferiamo utilizzare il sistema di Gauss in quanto fa intervenire soltanto la velocità della luce come costante fondamentale, anziché la permittività del vuoto $\epsilon_0$ e la permeabilità magnetica del vuoto $\mu_0$ (legate a $c$ dalla relazione $c=1/\sqrt{\epsilon_0\mu_0}$\,), con la conseguenza che tutte le formule risultano più semplici, ma soprattutto le equazioni fondamentali dei campi elettromagnetici nel vuoto esibiscono esplicitamente la simmetria fra il campo elettrico ed il campo magnetico, simmetria che ha carattere fondamentale in quanto richiesta, anzi imposta, dalla teoria della Relatività.\\
Tuttavia le unità pratiche, quelle con le quali abbiamo una certa familiarità: il volt, l'ampère, l'ohm e altre, sono unità del sistema MKSA e se non vogliamo perdere il contatto con i fatti più elementari dell'elettromagnetismo ed una certa sensibilità per gli ordini di grandezza, dovremo acquisire una certa familiarità con i fattori di ragguaglio fra le unità di misura nei due sistemi.

La (1.1) è una legge sperimentale, stabilita per corpi di piccole dimensioni o, equivalentemente, per distanze grandi rispetto alle dimensioni dei due corpi; nelle stesse condizioni si verifica sperimentalmente che la forza che due corpi di cariche $q_2$ e $q_3$, anche in posizioni diverse, esercitano su un corpo di carica $q_1$ è la somma delle forze $\vec F_{12}$ e $\vec F_{13}$ che ciascuno di essi eserciterebbe da solo sul corpo di carica $q_1$: in altri termini, le forze esercitate dalle cariche si sommano vettorialmente (\textbf{principio di sovrapposizione}). Tuttavia la (1.1) non ci è di grande aiuto se vogliamo determinare la forza che si esercita fra corpi di dimensioni non trascurabili, o anche nel caso in cui alcuni di essi siano conduttori, in quanto in tal caso, per il fenomeno dell'induzione elettrostatica, la disposizione delle cariche sui corpi dipende fortemente dalla distanza fra i corpi stessi. Allora assumiamo la (1.1) insieme al principio di sovrapposizione
come ``legge costitutiva'' valida per cariche puntiformi. Ciò significa che attribuiamo alla (1.1) un significato puramente matematico, ma da essa -- grazie al principio di sovrapposizione -- è possibile ricavare la forza che si esercita fra due qualsivoglia corpi carichi, di dimensioni arbitrarie: basta considerare ogni corpo carico come supporto di cariche puntiformi (opportunamente distribuite sul corpo), e calcolare la forza che si esercita su uno di essi come il risultante delle forze che ogni carica puntiforme di ogni altro corpo esercita sulle cariche puntiformi del corpo stesso. Dopo tutto è proprio quello che si fa per le interazioni gravitazionali: la legge di gravitazione di Newton si assume valida per masse puntiformi, e da essa si può ricavare la forza che si esercita fra corpi di dimensioni e forma arbitrarie. Il ruolo, nell'uno come nell'altro caso, del principio di sovrapposizione è fondamentale: anche se esprime un concetto, quello dell'additività vettoriale delle forze, che consideriamo scontato, è un principio del quale avremo molte occasioni di apprezzare l'importanza e la portata.\\
In ogni situazione sarà bene ricordarsi che la (1.1) per cariche puntiformi è una legge matematica: le cariche puntiformi in natura non esistono, e se dovessero sorgere dei problemi dovremo saper distinguere quali sono realmente problemi fisici e quali invece dovuti al concetto matematico di carica puntiforme (ambiguità, divergenze, energie infinite, ecc.).

\subsection{Il campo elettrico}

Sia data una carica puntiforme $q_1$ nella posizione $\vec r_1$. Definiamo \emph{campo elettrico} generato da $q_1$
\begin{equation}
\vec E_1(\vec r\,)\defeq k\,\frac{q_1}{|\vec r-\vec r_1|^3}\,(\vec r-\vec r_1)
\tag{1.4}
\end{equation}
cioè ad ogni punto di $R^3$, ad esclusione del punto $\vec r=\vec r_1$, è associato il vettore $\vec E_1(\vec r\,)$; la carica $q_1$ è detta la \emph{sorgente del campo}. Il campo elettrico ha le dimensioni di una forza diviso una carica e quindi si misura in dyne/franklin nel sistema di Gauss, e in newton/coulomb nel SI.\\
Il significato di $\vec E_1$ è il seguente: se nel punto (generico) $\vec r$ è presente una carica (puntiforme) $q$, su di essa si esercita una forza che per la (1.1) è data da
\begin{equation}
\vec F=q\vec E_1(\vec r\,)\,.
\tag{1.5}
\end{equation}
Se la sorgente è fissa nel punto $\vec r_1$, la (1.5) resta valida qualunque sia la carica $q$ nel punto $\vec r$; in tal caso $\vec F$ è proporzionale a $q$ ed il campo elettrico $\vec E_1(\vec r\,)$ è la forza che si esercita su una carica $q$ (in $\vec r\,$) divisa per $q$, ed è indipendente da $q$, e quindi -- in questo caso -- il campo elettrico è la forza agente sull'unità di carica. Se la sorgente non è fissa (per esempio è attaccata ad una molla) ma $\vec r_1$ è la sua posizione quando la carica $q$ si trova nel punto $\vec r$, è ancora vero \big(grazie alla (1.1)\big) che $q\vec E_1(\vec r\,)$ è la forza che $q_1$ esercita sulla carica $q$, ma non è più vero che detta forza è proporzionale a $q$, in quanto al variare di $q$ può cambiare la posizione della sorgente.\\
Se abbiamo $n$ cariche (puntiformi) $q_i$ nelle posizioni $\vec r_i$ il campo elettrico è definito (per ogni $\vec r\neq\vec r_i$) come la somma dei campi elettrici $\vec E_i(\vec r\,)$ generati dalle singole cariche $q_i$:
\begin{equation}
\vec E(\vec r\,)\defeq k\sum_{i=1}^n\frac{q_i}{|\vec r-\vec r_i|^3}\,(\vec r-\vec r_i)\,.
\tag{1.6}
\end{equation}
Si noti che le posizioni delle sorgenti sono punti di singolarità del campo elettrico.\\
Il significato di $\vec E(\vec r\,)$ è -- grazie al principio di sovrapposizione -- sempre quello discusso sopra: $\vec E(\vec r\,)$ è il rapporto fra la forza che si esercita su una carica $q$ posta in $\vec r$ e la carica stessa, quando le sorgenti $q_i$ sono nelle posizioni $\vec r_i$. In ogni caso la definizione di $\vec E(\vec r\,)$ \underline{prescinde} dal fatto che nel punto considerato sia presente o meno una carica, (la cosiddetta ``carica di prova'', o ``carica esploratrice'', o anche ``carica spia''), quindi quella che abbiamo dato non è una definizione ``operativa'', cioè che stabilisce quali misure occorre fare per determinare il campo. Sovente si definisce il campo elettrico (in un punto) come il rapporto fra la forza e la carica, quando questa sia così piccola da non influenzare le sorgenti: questa è solo apparentemente una definizione operativa, in quanto la carica più piccola che esiste in natura è quella dell'elettrone, per cui, per esempio, non sarebbe possibile definire il campo elettrico dentro un conduttore dove le sorgenti non fisse sono proprio gli elettroni, e quindi influenzabili da parte di un elettrone ``di prova''.\\
A questo punto è legittimo chiedersi se la definizione di campo elettrico sia di qualche utilità. In effetti, anche alla luce della discussione fatta sopra, può sembrare che il concetto di campo elettrico non sia particolarmente rilevante: se le sorgenti sono fisse, il campo elettrico permette sì di calcolare la forza che esse esercitano su una qualsiasi carica in un punto qualsiasi dello spazio, ma si tratta semplicemente di una riscrittura della legge di Coulomb più il principio di sovrapposizione, riscrittura -- tuttavia -- che evidenzia la proporzionalità della forza alla carica. Se le sorgenti non sono fisse, le loro posizioni sono determinate dalla posizione e dal valore di $q$ (oltre che da eventuali altre forze), e determinarle significa risolvere il problema delle posizioni di equilibrio di un sistema di $n+1$ punti materiali (le $q_i$ e $q$): risolto questo problema è già risolto anche quello di determinare la forza che le $q_i$ esercitano su $q$. Quindi non sembra che il concetto di campo elettrico sia di grande aiuto. Ma il fatto è che la domanda che ci siamo posti è intempestiva: per ora serve a farci capire che i problemi di interazione fra (molte) cariche elettriche sono problemi complicati. L'elettromagnetismo (l'elettrostatica è un caso particolare) affronta questi problemi da un punto di vista radicalmente diverso e innovativo: le incognite primarie del problema non sono più le forze, ma i campi (elettrico e magnetico); è quindi una teoria di campo, e non una teoria meccanica di molte particelle, e ha come obiettivi quello di stabilire le equazioni generali che devono essere soddisfatte dai campi (equazioni di Maxwell) e gli strumenti matematici per risolverle.\\
In questo primo capitolo ci limiteremo a studiare i casi in cui le sorgenti sono fisse: di ogni problema è allora possibile fornire la soluzione (la cui determinazione esplicita può comunque richiedere complicate tecniche di calcolo); la vera natura dell'elettrostatica come teoria di campo la incontreremo nel prossimo capitolo in cui tratteremo l'elettrostatica dei conduttori, in quanto sui conduttori le cariche non sono fisse, ma libere di muoversi sotto l'effetto delle loro interazioni reciproche e con eventuali altre cariche: tipico delle teorie di campo è questo duplice ruolo delle cariche sia come sorgenti del campo sia come oggetto di azioni da parte di esso.

Una utile rappresentazione grafica del campo elettrico si può ottenere mediante le \emph{linee di forza} del campo: sono linee orientate che in ogni punto sono tangenti al campo
elettrico in quel punto e hanno lo stesso verso. Per esempio, le linee di forza del campo generato da una sola carica puntiforme sono le semirette uscenti dalla posizione in cui si trova la sorgente, il verso è quello uscente dalla sorgente se questa ha carica positiva, entrante se ha carica negativa. Se il campo è generato da più cariche puntiformi, le linee di forza hanno in generale un andamento complicato; tuttavia, grazie alla singolarità $1/r^2$ del campo generato da ogni sorgente, tanto più ci si avvicina ad una sorgente, tanto meno il campo differisce da quello generato dalla sola sorgente vicina, come se le altre non ci fossero, quindi anche le linee di forza hanno in un intorno di ogni sorgente lo stesso andamento delle linee di forza di una carica isolata.\\
In tutti i punti in cui sono definite, le linee di forza sono curve regolari che non si intersecano; i punti singolari sono quelli in cui sono presenti le sorgenti, dove possono partire o arrivare infinite linee di forza, oppure quelli in cui $\vec E=0$ (p.es.\ il punto di mezzo della congiungente due cariche uguali), in un intorno dei quali il tipico andamento delle linee di forza è quello rappresentato in figura 1.1.
\fig{fig1_1}{1.1}
Ogni linea di forza che esce da una sorgente positiva o va all'infinito o, se non si arresta in un punto in cui $\vec E=0$, finisce in una sorgente negativa. Impareremo nel seguito maggiori dettagli sulle proprietà delle linee di forza; per esempio vedremo che la densità delle linee di forza in prossimità di un punto è proporzionale all'intensità del campo elettrico in quel punto: per fornire una anticipazione di questo risultato, consideriamo una sola sorgente (positiva) e $n$ linee di forza uscenti simmetricamente da essa; ogni sfera con centro nella sorgente (raggio $r$) è attraversata da tutte le $n$ linee di forza, la cui densità (numero/superficie) è quindi inversamente proporzionale alla superficie della sfera, cioè a $r^2$, come $E$. Quindi le linee di forza forniscono una buona rappresentazione non solo della direzione e del verso del campo elettrico ma, in virtù della legge $1/r^2$, anche dell'intensità.

\subsection{L'approssimazione del discreto con il continuo: distribuzioni di cariche}

Da un lato le nostre cariche puntiformi hanno solo carattere ``costitutivo'', cioè servono per costruire le cariche (estese) reali (persino l'elettrone è esteso: $r\approx10^{-12}\,\mathrm{cm}$), d'altro lato anche le distribuzioni di carica sui corpi macroscopici (bacchette di vetro, palline di sambuco, conduttori) sono costituite da numerosissime cariche di piccole dimensioni (elettroni, ioni) a distanze dell'ordine delle dimensioni atomiche ($\approx10^{-8}\,\mathrm{cm}$). Queste distribuzioni, discrete sulla scala microscopica, sono, alla scala macroscopica, convenientemente approssimate con distribuzioni continue. In modo analogo alla definizione di densità di massa, definiamo la densità di carica in un punto $P$ formalmente come
\begin{equation}
\rho(P)\defeq\lim_{\Delta V\to0}\frac{q}{\Delta V}
\tag{1.7}
\end{equation}
dove $\Delta V$ è un volumetto attorno al punto $P$, $q$ è la carica totale contenuta in esso, e $\Delta V\to0$ significa che la massima dimensione di $\Delta V$ tende a zero ($\Delta V$ tende a zero non solo come volume, ma anche come forma). In effetti il limite nella (1.7) fa parte solo della definizione formale: a causa della struttura (microscopicamente) discreta della distribuzione di carica, affinché $\rho(x,y,z)$ risulti una funzione continua (salvo la presenza di discontinuità su scala macroscopica come p.es.\ nel passaggio da una bacchetta elettrizzata al vuoto), è necessario che il processo di limite si arresti quando ancora $\Delta V$ contiene un gran numero di cariche microscopiche. In questo modo $\rho$ è definita in ogni punto come la densità media in un intorno -- macroscopicamente piccolo ma microscopicamente grande -- del punto stesso. Così, se consideriamo un corpo non elettrizzato (scarico) e non soggetto ad un campo elettrico esterno che ne alteri macroscopicamente la distribuzione microscopica delle cariche (come accade nei pezzettini di carta attirati dalla bacchetta di vetro elettrizzata o in un conduttore scarico vicino alla bacchetta di vetro), la densità macroscopica di carica $\rho$ risulta nulla: è bene rendersi conto che se invece veramente facessimo tendere a zero il $\Delta V$ nella (1.7), otterremmo una densità -- la densità microscopica -- che può variare, su distanze dell'ordine di $10^{-8}\,\mathrm{cm}$, da (circa) $+10^{22}\,\mathrm{C/m^3}$ (laddove è presente un protone) a $-10^{22}\,\mathrm{C/m^3}$ (dove c'è un elettrone), mentre una densità macroscopica dell'ordine di $10^3\,\mathrm{C/m^3}$ può già essere considerata una densità ragguardevole. È anche chiaro che con la definizione che abbiamo dato, in condizioni macroscopicamente stazionarie $\rho(x,y,z)$ risulta indipendente dal tempo, nonostante il moto di agitazione termica degli ioni o degli elettroni liberi (nei conduttori): se $\Delta V$ contiene un gran numero di cariche microscopiche, le variazioni nel tempo della carica in esso contenuta non sono macroscopicamente significative; si cerchi di immaginare quanto disordinatamente, invece, dipenderebbe dal tempo detta densità microscopica nel limite $\Delta V\to0\,$!\\
Contrariamente alla densità di massa, la densità di carica può assumere valori sia positivi che negativi.\\
Viceversa, se è nota la densità di carica $\rho(\vec r\,)$, è possibile calcolare la carica $Q_V$ contenuta in un arbitrario volume $V$. Il procedimento è del tutto simile a quello utilizzato per calcolare l'area sottesa dalla curva $y=f(x)$ e che porta alla definizione elementare di integrale definito della funzione $f(x)$ (integrale di Riemann), con l'unica differenza del numero di dimensioni, che qui sono tre: supponiamo che $\rho(\vec r\,)$ sia una funzione continua nel volume $V$ (insieme limitato e chiuso); sia $\Delta V_1\cdots\Delta V_i\cdots\Delta V_n$ una partizione di $V$ e $\vec r_i$ un punto arbitrario interno a $\Delta V_i$, allora quanto più $\Delta V_i$ è piccolo, tanto meglio $\rho(\vec r_i)\Delta V_i$ approssima la carica contenuta nell'i--esimo volumetto, per cui
\begin{equation}
Q_V=\lim_{n\to\infty}\sum_{i=1}^n\rho(\vec r_i)\Delta V_i\equiv\int_V\rho(\vec r\,)\,dV
\tag{1.8}
\end{equation}
dove il limite $n\to\infty$ viene realizzato facendo tendere a zero il massimo diametro dei volumetti $\Delta V_i$.\\
In modo analogo si procede per calcolare il campo elettrico generato da una distribuzione di carica $\rho$ contenuta in un volume $V$: come prima suddividiamo il volume $V$ in $n$ volumetti $\Delta V_i$. La carica contenuta nell'i--esimo volumetto è $q_i=\rho(\vec r_i)\Delta V_i$, dove $\vec r_i$ è un punto interno al volumetto (teorema della media), e contribuisce al campo elettrico nel punto $\vec r$
\begin{equation}
\vec E_i(\vec r\,)=k\,\frac{q_i}{|\vec r-\vec r_i|^3}\,(\vec r-\vec r_i)=k\,\frac{\rho(\vec r_i)}{|\vec r-\vec r_i|^3}\,(\vec r-\vec r_i)\,\Delta V_i\,.
\tag{1.9}
\end{equation}
La (1.9) vale se $|\vec r-\vec r_i|$ è molto maggiore delle dimensioni del volumetto $\Delta V_i$ quindi il campo elettrico totale, in qualsiasi punto $P$, si ottiene sommando i contributi $\vec E_i(\vec r\,)$
(principio di sovrapposizione) quando tutti i volumetti $\Delta V_i$ tendono a zero, cioè nel limite $n\to\infty$: in tal modo il punto $P$ finisce per essere (infinitamente) lontano da ogni volumetto, ad eccezione di quello che tende a zero avendo sempre $P$ come punto interno, ma vedremo che nel limite $\Delta V\to0$ questo porta contributo nullo al campo in $P$. Quindi
\begin{equation}
\begin{split}
\vec E(\vec r\,)&=\lim_{n\to\infty}\sum_{i=1}^n\vec E_i(\vec r\,)=\lim_{n\to\infty}k\sum_{i=1}^n\frac{\rho(\vec r_i)}{|\vec r-\vec r_i|^3}\,(\vec r-\vec r_i)\,\Delta V_i\\
&=k\int_V\frac{\rho(\vec r\,')}{|\vec r-\vec r\,'|^3}\,(\vec r-\vec r\,')\,dV.
\end{split}
\tag{1.10}
\end{equation}
Il campo elettrico prodotto dalla densità di carica macroscopica $\rho$ è esso stesso il campo elettrico macroscopico cioè, come $\rho$, la media del campo elettrico microscopico nell'intorno di ogni punto.\\
Nella (1.10) $\vec r$ è il punto in cui si calcola il campo elettrico ed $\vec r\,'$ il generico punto, all'interno del volume $V$, in cui è presente la sorgente $\rho(\vec r\,')$ e rispetto al quale viene effettuata l'integrazione. Dalla (1.10) può sembrare, a causa della presenza del denominatore $|\vec r-\vec r\,'|^3$, che il punto $\vec r$ debba essere esterno al volume $V$, così come nella (1.6) $\vec r$ deve essere diverso da ogni $\vec r_i$: no, la singolarità per $\vec r=\vec r\,'$ nella (1.10) è solo apparente, come vedremo fra poco, per cui $\vec E$ (se $\rho$ è ovunque finita) è definito anche dentro $V$, dove sono presenti le sorgenti.\\
Non avremo occasione di calcolare l'integrale di volume che compare nella (1.10) se non forse in qualche caso particolarmente semplice, tuttavia questo è il momento giusto per richiamare -- o definire -- coordinate diverse da quelle cartesiane ed il loro uso nel calcolo di integrali di volume. Il simbolo $dV$ che compare in un integrale di volume è il ricordo dei volumetti $\Delta V_i$ su cui si effettua la somma che porta alla definizione di integrale, e per questa ragione è detto ``elemento di volume infinitesimo'', e la sua espressione in termini delle coordinate dipende -- come quella dei volumetti finiti $\Delta V_i$ -- dalle coordinate utilizzate: se utilizziamo coordinate cartesiane, $dV=dx\,dy\,dz$; se utilizziamo coordinate cilindriche $(r,\theta,z)$, $dV=dr\cdot r d\theta\cdot dz=r\,dr\,d\theta\,dz$; se utilizziamo le \textbf{coordinate sferiche} o \textbf{coordinate polari} definite da (v.\ figura 1.2)
\fig{fig1_2}{1.2}
\begin{equation}
\left\{\begin{aligned}
x&=r\sen\theta\cos\phi\\
y&=r\sen\theta\sen\phi\\
z&=r\cos\theta
\end{aligned}\right.
\qquad\quad r\ge0;\quad 0\le\theta\le\pi;\quad 0\le\phi<2\pi
\tag{1.11}
\end{equation}
($r$ è la distanza dall'origine, $\theta$ e $\phi$ sono rispettivamente le coordinate angolari sulla superficie della sfera di raggio $r$, cioè la colatitudine e la longitudine), l'elemento di volume è dato da
\begin{equation}
dV=r^2\sen\theta\,dr\,d\theta\,d\phi\equiv r^2\,dr\,d\Omega\qquad\quad(d\Omega\equiv\sen\theta\,d\theta\,d\phi)
\tag{1.12}
\end{equation}
infatti $dS=r\sen\theta\,d\phi\cdot r\,d\theta=r^2\sen\theta\,d\theta\,d\phi$ è l'elemento di superficie sulla sfera di raggio $r$ fra due paralleli $\theta$, $\theta+d\theta$ e due meridiani $\phi$, $\phi+d\phi$ e $dV=dS\cdot dr$. Nella (1.12) $d\Omega$ è \emph{l'elemento di angolo solido}, cioè il rapporto fra $dS$ e $r^2$.

A questo punto possiamo verificare che la singolarità nella (1.10) è solo apparente: sia $P$ il punto (interno a $V$) nel quale vogliamo calcolare $\vec E$; prendiamo l'origine in $P$ e usiamo coordinate polari. Allora dalla (1.10) con $\vec r=0$ e dalla (1.12)
\begin{equation}
\begin{split}
\vec E(P)&=k\int_V\frac{\rho(\vec r\,')}{r'^2}\,\hat r\,'\,dV=k\int_V\frac{\rho(\vec r\,')}{r'^2}\,\hat r\,'\,r'^2\,dr'\,d\Omega\\
&=k\int_V\rho(\vec r\,')\,\hat r\,'\,dr'\,d\Omega
\end{split}
\tag{1.13}
\end{equation}
e la singolarità è sparita: intuitivamente si può dire che una distribuzione $\rho$ è costituita da cariche puntiformi infinitesime di ordine $\Delta V$, cioè $O(r^3)$, un ordine di infinitesimo maggiore dell'ordine di infinito della singolarità del campo di una carica puntiforme, che è $O(r^{-2})$.

La distribuzione $\rho$ è anche detta \emph{distribuzione di volume}; in modo analogo possiamo introdurre la \emph{distribuzione superficiale} $\sigma$, concetto particolarmente utile quando -- come nel caso dei conduttori -- le cariche sono distribuite in una regione di spessore molto piccolo, nei casi reali dell'ordine delle dimensioni atomiche:
\begin{equation}
\sigma(P)\defeq\lim_{\Delta S\to0}\frac{q}{\Delta S}
\tag{1.14}
\end{equation}
e già che ci siamo definiamo anche la \emph{distribuzione lineare} $\lambda$, che sarà utile quando p.es.\ le cariche sono distribuite su un sottile filo di quarzo:
\begin{equation}
\lambda(P)\defeq\lim_{\Delta l\to0}\frac{q}{\Delta l}\,\cdot
\tag{1.15}
\end{equation}
Il campo elettrico prodotto da una carica distribuita su una superficie $S$ similmente alla (1.10) è dato da
\begin{equation}
\vec E(\vec r\,)=k\int_S\frac{\sigma(\vec r\,')}{|\vec r-\vec r\,'|^3}\,(\vec r-\vec r\,')\,dS
\tag{1.16}
\end{equation}
e nel caso di una distribuzione lineare su una curva $\gamma$
\begin{equation}
\vec E(\vec r\,)=k\int_\gamma\frac{\lambda(\vec r\,')}{|\vec r-\vec r\,'|^3}\,(\vec r-\vec r\,')\,dl\,.
\tag{1.17}
\end{equation}
In presenza di una distribuzione superficiale il campo elettrico non è definito dove $\sigma\neq0$: l'argomento che ci ha portato alla (1.13) in questo caso non funziona perché l'elemento di superficie non ha l'$r^2\,dr$, ma solo $r\,dr$, non sufficiente a cancellare completamente la singolarità. Vedremo che ciò comporta soltanto che il campo elettrico subisce una discontinuità nell'attraversare una distribuzione superficiale, ma è finito da entrambe le parti della superficie. Nel caso invece di una distribuzione lineare il campo diverge come $1/r$, dove $r$ è la distanza dal filo. Quando queste singolarità ($1/r^2$ per le cariche puntiformi, $1/r$ per le distribuzioni lineari) o la discontinuità nel caso di distribuzioni superficiali, siano fonte di problemi o di ambiguità non di origine fisica, ma legate alla schematizzazione, questi problemi verranno risolti ricordando che distribuzioni superficiali, lineari e le cariche puntiformi, sono in realtà distribuzioni di volume in cui rispettivamente una, due o tre dimensioni sono state -- per comodità -- trascurate, e quindi vanno ripristinate.

\subsection{Il Flusso}

Sia data una superficie $S$ chiusa o no, e supponiamo che in ogni suo punto sia definita la normale alla superficie stessa: quindi $S$ è una superficie regolare, senza ``spigoli''. Supponiamo inoltre che $S$ sia \emph{orientabile}: ciò significa che, orientata (a piacere) la normale in un punto, per continuità risulta ben definito l'orientamento della normale in ogni altro punto di $S$. Un nastro di Moebius non è una superficie orientabile, in quanto dopo un giro si ritorna al punto di partenza con la normale invertita: in effetti le superfici possono essere estremamente complicate e presentare ogni tipo di patologia, ma noi avremo a che fare solo con superfici molto semplici come sfere, cilindri o, nei casi peggiori, parallelepipedi. È vero che i parallelepipedi hanno gli spigoli, ma sono comunque l'unione di un numero finito di superfici regolari, e ciò va ancora bene per i nostri scopi. Si dice che $S$ è ``orientata'' se ad ogni punto di $S$ è associato un versore normale alla superficie, variabile con continuità al variare del punto.\\
Sia $\vec F(\vec r\,)$ un campo vettoriale (un campo elettrico, oppure un campo magnetico, o il campo delle velocità di un fluido). Definiamo il \emph{flusso} $\Phi_S(\vec F)$ di $\vec F$ attraverso la superficie orientata $S$ nel modo seguente: in ogni punto di $S$ è definito il prodotto scalare $\vec F\cdot\hat n$, cioè la componente $F_n$ di $\vec F$ normale alla superficie; $\Phi_S(\vec F)$ è definito come l'integrale di $\vec F\cdot\hat n$ sulla superficie $S$:
\begin{equation}
\Phi_S(\vec F)\defeq\int_S\vec F\cdot\hat n\,dS\equiv\int_S\vec F\cdot d\vec S
\tag{1.18}
\end{equation}
dove nel secondo integrale abbiamo introdotto la definizione di elemento di area orientata $d\vec S\equiv\hat n\,dS$.\\
Il concetto di flusso, come altri concetti con cui avremo a che fare in seguito, trae le sue origini, e anche il nome, dalla teoria dei fluidi, e ad essa facciamo riferimento per capire il significato del flusso: supponiamo di avere un fluido in moto in un condotto, e supponiamo che il campo delle velocità del fluido sia costante: $\vec v(x,y,z)=\vec v_0$. Presa una superficie piana di area $A$, ortogonale a $\vec v_0$ e orientata come $\vec v_0$ (figura 1.3), dalla (1.18) si ha che il flusso di $\vec v$ attraverso $A$ è dato da
\fig{fig1_3}{1.3}
\begin{equation}
\Phi_A(\vec v)=v_0\,A
\tag{1.19}
\end{equation}
e siccome $v_0\Delta t\,A$ è il volume di fluido che attraversa la superficie $A$ nell'intervallo di tempo $\Delta t$, ne segue che $\Phi_A(\vec v)$ coincide con il volume di fluido che attraversa $A$ nell'unità di tempo, cioè quella che normalmente si chiama la \emph{portata} del fluido attraverso $A$. Se $A$ non è ortogonale a $\vec v_0$, ma la normale $\hat n$ ad essa forma un angolo $\theta$ con $\vec v_0$, chiaramente la portata non è $v_0\,A$, ma $v_0\,A\cos\theta$, cioè ancora il flusso attraverso $A$.

Se $\vec v$ non è costante su $A$, e/o $A$ non è piana, $\Phi_A(\vec v)$ è ottenuto come somma dei flussi attraverso una partizione di $A$ in $n$ superfici $\Delta S_i$ sufficientemente piccole da poter essere considerate piane, e da poter considerare $\vec v$ costante su ciascuna di esse, ed infine prendendo il limite $n\to\infty$ \big(la definizione dell'integrale di superficie della (1.18)\big).

Sia ora $S$ una superficie chiusa e $V$ il volume racchiuso da essa: in questo caso conveniamo che la normale $\hat n$ alla superficie sia (ovunque) orientata verso l'esterno, e quindi
parleremo di \emph{flusso uscente da $S$} o anche \emph{flusso uscente dal volume $V$}.\\
Consideriamo il campo elettrico generato da una carica puntiforme $q$ e calcoliamo il flusso uscente da una sfera di raggio $r$ con centro nella carica $q$; siccome in questo caso $\hat n=\hat r$, si ha
\begin{equation}
\Phi(\vec E)=\int k\,\frac{q}{r^2}\,\hat r\cdot\hat n\,dS=k\,\frac{q}{r^2}\int dS=k\,\frac{q}{r^2}\,4\pi r^2=4\pi kq.
\tag{1.20}
\end{equation}
Il risultato è indipendente dal raggio della sfera: si tratta evidentemente di una conseguenza della legge $1/r^2$. In effetti la (1.20) è un caso particolare del teorema di Gauss che dimostreremo nel prossimo paragrafo, ma del quale cercheremo ora di dare una versione che aiuti l'intuizione: si supponga che la carica $q$ sia una sorgente che emette, in tutte le direzioni e a ritmo costante, un fluido (gas, liquido, luce \dots) in moto stazionario (cioè ogni grandezza -- velocità, densità ecc.\ -- è in ogni punto indipendente dal tempo). Il flusso attraverso una superficie $S$ è la quantità di fluido che la attraversa nell'unità di tempo: è chiaro allora che il flusso -- cioè la quantità di fluido -- che attraversa una sfera con centro nella sorgente è indipendente dal raggio della sfera ed è uguale alla portata della sorgente: se così non fosse si avrebbe in qualche regione dello spazio un accumulo o una diminuzione di fluido, contro l'ipotesi di moto stazionario (quindi il nostro fluido è un buon modello per il campo elettrico). Ma allora è anche chiaro che detto flusso è indipendente dalla forma della superficie chiusa che contiene la sorgente, e che se la superficie (chiusa) non contiene la sorgente, il flusso totale uscente da essa è nullo: tanto fluido entra e altrettanto esce. Se poi abbiamo tante sorgenti, e se vale il principio di sovrapposizione, cioè i fluidi emessi da ciascuna di esse non interagiscono (gas ideali, fotoni, \dots), il flusso uscente da una superficie chiusa è uguale alla somma dei flussi delle sole sorgenti contenute all'interno, cioè la somma delle loro portate: e questo è proprio quello che dice il teorema di Gauss.

\subsection{Il teorema della divergenza e il teorema di Gauss}

È dato il solito campo vettoriale $\vec F(\vec r\,)$. Sia $P$ un punto, $\Delta V$ un volume che lo contiene e $\Phi_{\Delta V}$ il flusso uscente da $\Delta V$. Si definisce ``\emph{divergenza di $\vec F$}'' nel punto $P$ il limite -- quando esiste finito -- del rapporto flusso/volume al tendere a zero del volume attorno al punto $P$:
\begin{equation}
\dive\vec F\defeq\lim_{\Delta V\to0}\frac{\Phi_{\Delta V}(\vec F)}{\Delta V}\,.
\tag{1.21}
\end{equation}
La divergenza di un campo vettoriale è quindi la densità del flusso di volume ed è una funzione scalare associata al campo vettoriale $\vec F$: nei punti in cui $\dive\vec F>0$ si hanno le ``sorgenti del flusso'', mentre dove $\dive\vec F<0$ si hanno le sorgenti negative, cioè i ``pozzi''.\\
Abbiamo detto, a proposito del limite, ``quando esiste'': per esempio, nei punti in cui è presente una carica puntiforme la divergenza del campo elettrico non esiste in quanto, se $\Delta V$ è una sferetta con centro in $P$, $\Phi_{\Delta V}$ resta costante (è indipendente dal raggio) mentre $\Delta V$ tende a zero. Vedremo che la divergenza esiste in tutti i punti in cui $\vec F$ è derivabile (funzione di classe $C^1$).\\
Sia ora $V$ un volume arbitrario dentro il quale $\dive\vec F$ è ovunque definita e sia $\Phi_V$ il
flusso di $\vec F$ uscente da $V$. Se $V$ viene diviso in due parti $V_1$ e $V_2$ da una superficie $S'$ (figura 1.4), si ha
\begin{equation}
\Phi_V\equiv\Phi_{V_1\cup V_2}=\Phi_{V_1}+\Phi_{V_2}
\tag{1.22}
\end{equation}
\fig{fig1_4}{1.4}
infatti il flusso attraverso $S'$ ha un valore quando $S'$ è considerata come parte di $V_1$, e il valore opposto (cambia verso la normale $\hat n$) quando è considerata come parte di $V_2$. Quindi il flusso è una funzione additiva del volume, proprio come la massa o la carica, e la divergenza è una vera densità: se $\Delta V_1\cdots\Delta V_i\cdots\Delta V_n$ è una partizione di $V$, si ha
\begin{equation}
\Phi_V=\sum_{i=1}^n\Phi_{\Delta V_i}
\tag{1.23}
\end{equation}
quindi, come la carica (o la massa) contenuta in un volume $V$ è data dall'integrale della rispettiva densità, così
\begin{equation}
\Phi_V(\vec F)\equiv\int_S\vec F\cdot d\vec S=\int_V\dive\vec F\,dV\,.
\tag{1.24}
\end{equation}
La (1.24) costituisce il

\textbf{\textit{Teorema della divergenza.}} \textit{Il flusso di un campo vettoriale attraverso una superficie chiusa $S$ è uguale all'integrale della sua divergenza nel volume racchiuso da $S$.}

Il teorema è completo quando siamo in grado di calcolare direttamente $\dive\vec F$, noto il campo vettoriale $\vec F(\vec r\,)$: sia $P$ un punto (in cui $\dive\vec F$ esiste) e prendiamo per $\Delta V$ un parallelepipedo con centro in $P$. Siccome (per definizione) $\dive\vec F$ è indipendente dall'orientamento degli assi cartesiani, prendiamo gli assi coordinati paralleli agli spigoli del parallelepipedo; siano $x,y,z$ le coordinate di $P$ e $\Delta x,\Delta y,\Delta z$ le dimensioni del parallelepipedo. Calcoliamo, a meno di infinitesimi di ordine superiore in $\Delta x,\Delta y,\Delta z$ (che alla fine dovranno tendere a zero), il flusso di $\vec F$ uscente dal parallelepipedo come somma dei flussi attraverso le sue basi: il flusso $\Phi_z$ uscente dalle due basi ortogonali all'asse $z$ è
\begin{equation}
\begin{split}
\Phi_z&\simeq F_z\Big(x,y,z+\frac{\Delta z}{2}\Big)\Delta x\Delta y-F_z\Big(x,y,z-\frac{\Delta z}{2}\Big)\Delta x\Delta y\simeq\\
&\simeq\frac{\partial F_z}{\partial z}\,\Delta x\Delta y\Delta z
\end{split}
\tag{1.25}
\end{equation}
e analogamente per il flusso uscente dalle due basi ortogonali all'asse $y$ e per quello uscente dalle basi ortogonali all'asse $x$, per cui
\begin{equation}
\Phi_{\Delta V}=\Big(\frac{\partial F_x}{\partial x}+\frac{\partial F_y}{\partial y}+\frac{\partial F_z}{\partial z}\Big)\,\Delta x\Delta y\Delta z+o(\Delta x\Delta y\Delta z)
\tag{1.26}
\end{equation}
e quindi per la (1.21) l'espressione di $\dive\vec F$ in coordinate cartesiane è
\begin{equation}
\dive\vec F(x,y,z)=\frac{\partial F_x}{\partial x}+\frac{\partial F_y}{\partial y}+\frac{\partial F_z}{\partial z}\,\cdot
\tag{1.27}
\end{equation}

Calcoliamo la divergenza del campo elettrico di una carica puntiforme, in tutti i punti ad eccezione di quello occupato dalla carica. Presa l'origine delle coordinate nella posizione della carica, si ha
\begin{equation}
\left\{\begin{aligned}
E_x(x,y,z)&=kq\,\frac{x}{(x^2+y^2+z^2)^{3/2}}\\
E_y(x,y,z)&=kq\,\frac{y}{(x^2+y^2+z^2)^{3/2}}\\
E_z(x,y,z)&=kq\,\frac{z}{(x^2+y^2+z^2)^{3/2}}
\end{aligned}\right.
\tag{1.28}
\end{equation}
e, per $\vec r\neq0$,
\begin{equation}
\frac{\partial E_x}{\partial x}=kq\Big(\frac{1}{(x^2+y^2+z^2)^{3/2}}-\frac{3x^2}{(x^2+y^2+z^2)^{5/2}}\Big)
\tag{1.29}
\end{equation}
e analoghe, per cui, sommando
\begin{equation}
\dive\vec E(x,y,z)=\frac{\partial E_x}{\partial x}+\frac{\partial E_y}{\partial y}+\frac{\partial E_z}{\partial z}=0\quad\forall\vec r\neq0
\tag{1.30}
\end{equation}
quindi la divergenza del campo di una carica puntiforme è nulla in tutti i punti ad eccezione di quello in cui è presente la carica.

Abbiamo finalmente tutti gli strumenti per dimostrare l'annunciato teorema di Gauss. Sia $q$ una carica puntiforme ed $S$ una arbitraria superficie chiusa non contenente $q$ al suo interno. Per il teorema della divergenza (1.24) e per la (1.30), il flusso del campo elettrico attraverso $S$ è nullo.\\
Sia ora $S$ una superficie chiusa che contiene $q$ al suo interno: allora esiste una sfera con centro nella posizione di $q$ e tutta contenuta dentro $S$; il volume $V$ racchiuso da $S$ risulta diviso in due volumi (v.\ figura 1.5): la sfera $V_0$ ed il resto $V_1\equiv V\setminus V_0$, ed il flusso uscente da $V$ è la somma dei flussi uscenti dai due volumi
\fig{fig1_5}{1.5}
\begin{equation}
\Phi_V(\vec E)=\Phi_{V_0}(\vec E)+\Phi_{V_1}(\vec E)\,.
\tag{1.31}
\end{equation}
Il flusso $\Phi_{V_0}$ uscente dalla sfera è, per la (1.20), $4\pi kq$, mentre quello uscente dal volume rimanente $V_1$ è nullo, in quanto non contiene la carica. Quindi il flusso uscente da una superficie chiusa è nullo se non contiene la carica, e vale $4\pi kq$ se contiene la carica. Se abbiamo più cariche, siccome il campo elettrico è la somma dei campi elettrici ed il flusso, essendo lineare nel campo, è la somma dei flussi, il teorema dimostrato per una carica resta valido per un numero qualsiasi di cariche, e quindi anche nel caso limite di distribuzioni continue:

\textbf{\textit{Teorema di Gauss.}} \textit{Il flusso del campo elettrico uscente da una superficie chiusa è proporzionale alla carica totale contenuta dentro la superficie:}
\begin{equation}
\Phi_S(\vec E)=4\pi k\,Q^{\rm int}.
\tag{1.32}
\end{equation}
E se una carica puntiforme si trova proprio sulla superficie $S$? Ecco uno di quei casi ``ambigui'', citati alla fine del 1$^\circ$ paragrafo, imputabile alla schematizzazione che
fa ricorso alle cariche puntiformi: nei casi reali ogni carica ha dimensioni non nulle, (p.es.\ la nostra carica è distribuita dentro una sferetta di raggio finito), e quella che contribuisce a 2$^\circ$ membro della (1.32) è solo la porzione della carica (estesa) interna al volume $V$.

Il teorema di Gauss è conseguenza della legge di Coulomb, cioè della legge $1/r^2$ e, con l'ipotesi che il campo di una carica puntiforme (in quiete) è a simmetria sferica (campo centrale), è ad essa equivalente: infatti, se $\vec E(\vec r\,)=E_r(r)\,\hat r$, presa una sfera di raggio $r$ con centro nella carica $q$, dalla (1.32) si ha
\begin{equation}
4\pi kq=\Phi_S(\vec E)=\int_S E_r(r)\,dS=4\pi r^2E_r(r)\quad\Rightarrow\quad\vec E(\vec r\,)=k\,\frac{q}{r^2}\,\hat r
\tag{1.33}
\end{equation}
quindi la (1.32), più l'ipotesi di simmetria sferica, assume il ruolo di equazione fondamentale dell'elettrostatica. Essa è una equazione di carattere ``integrale'', cioè concerne regioni finite di spazio, e può essere trasformata in una equazione di carattere ``locale'', cioè valida punto per punto: sia $\Delta V$ un volume arbitrario contenente il punto $P$ in cui è presente solo una distribuzione di volume $\rho$, e sia $q$ la carica in esso contenuta; per la (1.21) (definizione di divergenza) e la (1.32) (teorema di Gauss) si ha
\begin{equation}
\dive\vec E\equiv\lim_{\Delta V\to0}\frac{\Phi_{\Delta V}(\vec E)}{\Delta V}=4\pi k\lim_{\Delta V\to0}\frac{q}{\Delta V}\equiv4\pi k\rho(P)
\tag{1.34}
\end{equation}
cioè
\begin{equation}
\dive\vec E(x,y,z)=4\pi k\rho(x,y,z)
\tag{1.35}
\end{equation}
che è, appunto, la versione locale del teorema di Gauss. La (1.35), tuttavia, è valida solo nei punti in cui si ha (eventualmente) una distribuzione di volume: sono escluse le distribuzioni superficiali o lineari, oltreché le cariche puntiformi, quindi in questo senso la (1.32) è più generale della (1.35). Tuttavia, se ci ricordiamo che tutte le citate distribuzioni ``singolari'' debbono essere considerate come limiti di distribuzioni di volume, riconosciamo il carattere generale della (1.35), che quindi costituisce la prima equazione fondamentale dell'elettromagnetismo, o \textbf{\textit{1$^{\mathbf{a}}$ equazione di Maxwell.}}

Concludiamo questo paragrafo dimostrando la ``legge di densità'' delle linee di forza, citata alla fine del 2$^\circ$ paragrafo: in una regione priva di sorgenti consideriamo un punto $P_0$ e una superficie ``infinitesima'' $\Delta S_0$ contenente $P_0$ e ortogonale ad $\vec E(P_0)$; l'insieme delle linee di forza che attraversano $\Delta S_0$ costituiscono quello che si chiama un \emph{tubo di flusso} di sezione infinitesima (figura 1.6); sezioniamo questo tubo di flusso con una superficie $\Delta S$ attorno ad un punto $P$, anch'essa ortogonale ad $\vec E(P)$. Il flusso di $\vec E$ uscente dalla porzione di tubo di flusso delimitata dalle due superfici $\Delta S_0$ e $\Delta S$ è nullo perché per ipotesi al suo interno non ci sono sorgenti; il flusso attraverso la sua superficie laterale è nullo perché, essendo costituita da linee di forza, $\vec E$ è tangente alla superficie, quindi il flusso entrante attraverso $\Delta S_0$ e quello uscente attraverso $\Delta S$ sono uguali:
\fig{fig1_6}{1.6}
\begin{equation}
E(P_0)\Delta S_0=E(P)\Delta S\quad\Rightarrow\quad\frac{E(P)}{E(P_0)}=\frac{\Delta S_0}{\Delta S}
\tag{1.36}
\end{equation}
cioè l'intensità del campo è inversamente proporzionale alla sezione del tubo di flusso. Se da $\Delta S_0$ facciamo partire $n$ linee di forza, queste stesse linee di forza, e solo esse, attraversano $\Delta S$, quindi anche la loro densità è inversamente proporzionale alla sezione del tubo di flusso, e quindi è proporzionale all'intensità del campo elettrico.

\subsection{Le leggi di trasformazione del campo elettrico}

La (1.6) nel caso di cariche puntiformi, le (1.10), (1.16) e (1.17) nel caso di distribuzioni di cariche, permettono -- in linea di principio -- di determinare il campo elettrico per qualunque distribuzione delle sorgenti. Ad eccezione della (1.6) (nel caso di poche cariche), le altre espressioni comportano anche nei casi più semplici calcoli alquanto complessi. Tuttavia, tutte le volte che le sorgenti presentino sufficienti proprietà di simmetria, lo strumento giusto per determinare il campo elettrico è rappresentato dal teorema di Gauss, e va anche detto che i problemi risolubili solo con la (1.10) (o le altre) e non mediante il teorema di Gauss, sono veramente pochi e si trovano trattati solo su testi di carattere avanzato.\\
Per poter utilizzare le semplificazioni derivanti dalle proprietà di simmetria del problema, è necessario conoscere le leggi di trasformazione del campo elettrico per particolari trasformazioni delle sorgenti.\\
Per semplicità di enunciato faremo riferimento al caso in cui le sorgenti sono descritte da una distribuzione di volume $\rho$, ma sarà ovvio che tutto quanto diremo vale tale quale per distribuzioni superficiali, lineari, e anche per cariche puntiformi. Cominciamo con il caso delle isometrie (useremo notazioni più adatte al tipo di problema che vogliamo discutere).\\
Sia data una distribuzione $\rho(x)$, $x\equiv(x_1,x_2,x_3)$, e sia $\vec E(x)$ il campo elettrico generato da $\rho$. Sia $\psi(x)$ una isometria dello spazio:
\begin{equation}
\begin{gathered}
\psi(x)\equiv Rx+a;\quad\psi^{-1}(x)=R^{-1}(x-a);\quad\Big(\psi(x)_i=\sum_{j=1}^3R_{ij}x_j+a_i\Big)\\
{}^t\!RR=I;\ \det R=\pm1
\end{gathered}
\tag{1.37}
\end{equation}
e consideriamo la distribuzione
\begin{equation}
\rho'(x)=\rho\big(\psi^{-1}(x)\big)=\rho\big(R^{-1}(x-a)\big).
\tag{1.38}
\end{equation}
Il significato di $\rho'(x)$ è questo: siccome $\rho'\big(\psi(x)\big)=\rho(x)$, la nuova distribuzione $\rho'$ nel punto trasformato $(\psi(x))$ è uguale alla vecchia distribuzione $\rho$ nel punto non trasformato $(x)$, cioè $\rho'$ si ottiene assoggettando rigidamente la distribuzione $\rho$ alla stessa trasformazione a cui l'isometria $\psi$ assoggetta i punti dello spazio: $x\to\psi(x)$ (questo è quello che si chiama ``punto di vista attivo''). Così, per esempio, nel caso della traslazione $\psi(x)=x+a$, se $\rho$ è concentrata attorno al punto $x_0$, $\rho'$ è concentrata attorno al punto $x_0+a$. Nella figura 1.7, nella quale oltre alle distribuzioni $\rho$ e $\rho'$ sono rappresentati anche i campi elettrici generati da ciascuna di esse, l'isometria consiste in una rotazione di 90$^\circ$ in verso orario attorno al punto $P$, seguita dalla traslazione individuata dal vettore $\vec a\equiv\overline{PP'}$.
\fig{fig1_7}{1.7}
Come illustrato dalla figura 1.7, è evidente (e lo dimostreremo) che se $\rho'$ si ottiene traslando, ruotando o riflettendo $\rho$ in un certo modo, il campo elettrico $\vec E\,'(x)$ generato dalla distribuzione $\rho'(x)$ si ottiene traslando, ruotando o riflettendo nello stesso modo il campo $\vec E(x)$, generato dalla distribuzione $\rho(x)$: in altri termini, la distribuzione ``si porta dietro il proprio campo''.\\
Quanto appena affermato è espresso dal

\textbf{\textit{Teorema.}} \textit{Il campo elettrico $\vec E'(x)$ generato dalla distribuzione $\rho'(x)$ è dato da}
\begin{equation}
E'_i(x)=\textstyle\sum_j R_{ij}E_j\big(\psi^{-1}(x)\big).
\tag{1.39}
\end{equation}
Se nella (1.39) poniamo $\psi(x)$ al posto di $x$ si ottiene, analogamente alla (1.38), $E'_i\big(\psi(x)\big)=\sum_jR_{ij}E_j(x)$; la presenza della matrice di rotazione $R$ a secondo membro della (1.39) e non della (1.38) è dovuta al fatto che $\vec E$ è un campo vettoriale, mentre $\rho$ è una funzione (campo) scalare: ruotare $\vec E$ significa trasportarlo nel punto ruotato, ruotandone conformemente la direzione.\\
La (1.39), è la legge generale di trasformazione di un qualsiasi campo vettoriale (campi di forze, campi di velocità, ecc.); per i campi pseudovettoriali (vettori assiali), come il campo magnetico, la legge di trasformazione differisce dalla (1.39), nel caso di trasformazioni con $\det R=-1$, per il segno.\\
Se dimostriamo il teorema per una carica puntiforme, il principio di sovrapposizione ci garantisce che vale per un numero qualsiasi di cariche e quindi anche per una distribuzione continua. Sia $x_0$ la posizione della carica $q$. L'isometria $\psi$ porta la carica nel punto $\psi(x_0)=Rx_0+a$, quindi
\begin{equation}
\begin{split}
E'_i(x)&=kq\,\frac{x_i-(Rx_0+a)_i}{\|x-(Rx_0+a)\|^3}=kq\,\frac{(x-a-Rx_0)_i}{\|x-a-Rx_0\|^3}\\
&=\textstyle\sum_jR_{ij}\,kq\,\dfrac{\big(R^{-1}(x-a)-x_0\big)_j}{\|R^{-1}(x-a)-x_0\|^3}=\sum_jR_{ij}E_j\big(\psi^{-1}(x)\big)\quad\blacksquare
\end{split}
\tag{1.40}
\end{equation}
Nella dimostrazione (penultimo passaggio) abbiamo solo utilizzato il fatto che la distanza $\|x-x'\|$ fra due punti è invariante per isometrie, e non la particolare dipendenza di $\vec E$ dalla distanza, quindi la (1.39) è indipendente dalla legge $1/r^2$.

Casi particolari:
\begin{equation}
\begin{aligned}
&\psi(x)=x+a&\qquad&\vec E'(\vec r\,)=\vec E(\vec r-\vec a)\\
&\psi(x)=Rx&&E'_i(x)=\textstyle\sum_jR_{ij}E_j(R^{-1}x)\\
&\psi(x)=-x&&\vec E'(\vec r\,)=-\vec E(-\vec r\,).
\end{aligned}
\tag{1.41}
\end{equation}
\textbf{\textit{Corollario.}} \textit{Se} $\rho\big(\psi(x)\big)=\rho(x)\quad\Rightarrow$
\begin{equation}
E_i\big(\psi(x)\big)=\textstyle\sum_jR_{ij}E_j(x).
\tag{1.42}
\end{equation}
Infatti $\rho(x)=\rho\big(\psi(x)\big)\ \Leftrightarrow\ \rho'(x)=\rho(x)$ e siccome il campo elettrico è univocamente determinato dalle sorgenti, $\vec E\,'(\vec r\,)=\vec E(\vec r\,)$ quindi dalla (1.39) con $\psi(x)$ al posto di $x$ si ha subito la tesi.\\
Anche il significato di questo corollario è evidente: se le sorgenti sono invarianti per una
certa trasformazione (isometrica), anche il campo elettrico globalmente resta invariato per la stessa trasformazione: p.es.\ se le sorgenti sono invarianti per una rotazione (o una traslazione), l'intero campo elettrico (pensato come un corpo rigido) resta invariato per la stessa rotazione (o traslazione).\\
Tutto sommato, sia il teorema che il corollario esprimono dei concetti abbastanza intuitivi, tanto che il loro tacito uso è implicito nella frase ``per ovvie ragioni di simmetria'', spesso utilizzata in molti testi. Concetti intuitivi, sì, ma non ``scontati'': alla base di tutto c'è il fatto che la legge di Coulomb è la stessa in qualsiasi luogo (lo spazio è ``omogeneo e isotropo''); se lo spazio in cui viviamo fosse pervaso da un fluido (\emph{etere}) non omogeneo e con proprietà dipendenti dalla direzione, cioè non isotropo, le ovvie ragioni di simmetria verrebbero meno.

L'insieme delle trasformazioni $\psi$ tali che $\rho\big(\psi(x)\big)=\rho(x)$ costituiscono chiaramente un gruppo (sottogruppo delle isometrie), che chiameremo ``il gruppo $\mathcal G$ di invarianza delle sorgenti''; è evidente che tanto più il gruppo di invarianza è ricco, e tante più informazioni si ottengono sul campo elettrico. Queste informazioni consistono in relazioni fra il campo elettrico nei vari punti appartenenti ad una stessa orbita per l'azione di $\mathcal G$ sui punti dello spazio (un'orbita è l'insieme dei punti che si ottiene da un dato punto per l'azione del gruppo su di esso): la (1.42) dice che la conoscenza di $\vec E$ in un punto $P_0$ determina il campo in tutti i punti $P=\psi(P_0)$ dell'orbita che contiene $P_0$; in particolare l'intensità di $\vec E$ è la stessa in tutti i punti di una stessa orbita.\\
Preso un punto $P$, lo stabilizzatore $\mathcal G_P$ di $P$ è il sottogruppo di $\mathcal G$ che lascia invariato il punto $P$; se $\mathcal G_P$ non è banale, si ottengono informazioni direttamente sulla direzione di $\vec E(P)$ (in generale non sul verso e neppure sull'intensità: la (1.42) è un'equazione lineare e omogenea), in quanto $\vec E(P)$ deve essere invariante (cioè autovettore) per tutte le trasformazioni di $\mathcal G_P$:
\begin{equation}
E_i(P)=\textstyle\sum_jR_{ij}E_j(P)\qquad\forall\,R\in\mathcal G_P.
\tag{1.43}
\end{equation}
$\mathcal G_P$ può consistere solo di rotazioni e/o riflessioni (infatti le traslazioni non lasciano fisso alcun punto); spesso (dipende da $\mathcal G_P$) queste informazioni determinano completamente la direzione di $\vec E$ nel punto $P$, e quindi in tutti i punti dell'orbita di $P$.\\
(Di solito il gruppo di invarianza viene caratterizzato elencando alcuni suoi sottogruppi che lo generano completamente, per esempio le rotazioni e le traslazioni, oppure le rotazioni e le riflessioni, ecc.).

\paragraph{Esempi}
\begin{enumerate}
\item Consideriamo il caso particolarmente importante di una distribuzione a simmetria sferica: $\rho=\rho(r)$. Il gruppo di invarianza delle sorgenti è il gruppo $O_3$ delle rotazioni e riflessioni; non sarà necessario utilizzarlo tutto: sono sufficienti le rotazioni. Sia $P$ un punto qualsiasi: se $P$ coincide con il centro $O$ della distribuzione, $\mathcal G_P$ è l'intero gruppo per cui, se $\vec E(O)$ è definito, per la (1.43) deve essere $E_i(O)=\sum_jR_{ij}E_j(O)$, e ciò è possibile solo se $\vec E(O)=0$ (se p.es.\ è presente una carica puntiforme nel centro, $\vec E(O)$ non è definito); altrimenti ($P\neq O$) $\mathcal G_P$ è il gruppo delle rotazioni attorno alla retta $OP$ ed $\vec E(P)$ è invariante solo se parallelo a detta retta, cioè se è radiale. L'orbita di $P$ è la superficie della sfera di centro $O$
passante per $P$ sulla quale quindi $\vec E$ ha intensità costante, è radiale ed è ovunque diretto o verso l'esterno o verso l'interno: in definitiva $\vec E(\vec r\,)=E_r(r)\,\hat r$.
\item Distribuzione con simmetria cilindrica; preso l'asse $z$ come asse di simmetria, $\rho=\rho(x^2+y^2)$.\\
Per esempio, un cilindro con carica uniformemente distribuita al suo interno, oppure un filo rettilineo indefinito con densità lineare $\lambda$ costante. Il gruppo di invarianza contiene le rotazioni attorno all'asse $z$, le traslazioni parallele all'asse $z$ e le rotazioni di 180$^\circ$ attorno ad ogni retta perpendicolare all'asse $z$. Ci sono anche le riflessioni, ma non ne avremo bisogno. Sull'asse $z$ $\vec E=0$ (se è definito); se $P$ è un punto non appartenente all'asse $z$, $\mathcal G_P$ è la rotazione di 180$^\circ$ attorno alla retta per $P$ ortogonale all'asse $z$, quindi $\vec E(P)$ deve essere parallelo a detta retta: campo ``radiale piano''; l'orbita di $P$ è il cilindro per $P$ avente per asse l'asse $z$. Se chiamiamo $\vec\varrho$ il vettore di componenti $(x,y,0)$, si ha $\vec E(\vec r\,)=E_\varrho(\varrho)\,\hat\varrho$.
\item Distribuzione con simmetria piana: $\rho=\rho(z)$. Per esempio, una carica uniformemente distribuita fra due piani paralleli, oppure un piano con densità superficiale costante $\sigma$. Il gruppo di invarianza contiene le traslazioni parallele al piano $z=0$ e le rotazioni attorno ad ogni asse parallelo all'asse z (gli esempi dati sopra hanno qualche simmetria in più). Le rotazioni attorno alla retta parallela all'asse $z$ passante per il punto $P$ costituiscono lo stabilizzatore di $P$, per cui $\vec E$ è ovunque parallelo all'asse $z$. I piani $z=$ costante sono le orbite sulle quali il campo è costante. Quindi $\vec E(\vec r\,)=\big(0,0,E_z(z)\big)$.
\end{enumerate}
Infine, talvolta ci sarà utile un'altra proprietà del campo elettrico: se cambiamo segno a tutte le sorgenti, il campo elettrico cambia segno:
\begin{equation}
\rho'(\vec r\,)=-\rho(\vec r\,)\quad\Rightarrow\quad\vec E'(\vec r\,)=-\vec E(\vec r\,).
\tag{1.44}
\end{equation}
La trasformazione $q_i\to-q_i$ (o $\rho(x)\to-\rho(x)$) è detta \emph{Coniugazione di carica}.\\
La coniugazione di carica non può essere una trasformazione di invarianza (salvo il caso banale $\vec E\equiv0$), ma può esserlo composta con altre trasformazioni discrete.\\
Per esempio, se abbiamo due distribuzioni simmetriche rispetto ad un piano (piano $z=0$) ma di cariche opposte, come in figura 1.8 (o anche nel caso di un \textbf{dipolo}: sistema costituito da due cariche $q$ e $-q$), la composizione di una riflessione rispetto al piano $z=0$ con la coniugazione di carica, $\rho'(x,y,z)\equiv-\rho(x,y,-z)$, è una trasformazione di invarianza delle sorgenti:
\fig{fig1_8}{1.8}
$\rho'(x,y,z)\equiv-\rho(x,y,-z)=\rho(x,y,z)$,\quad per cui dalle (1.39) e (1.44), se $R_{ij}$ è la matrice diagonale con gli elementi diagonali $R_{11}=R_{22}=1$, $R_{33}=-1$,
\begin{equation}
E_i(x,y,-z)=-R_{ii}E_i(x,y,z)
\tag{1.45}
\end{equation}
cioè
\begin{equation}
\left\{\begin{aligned}
E_x(x,y,z)&=-E_x(x,y,-z)\\
E_y(x,y,z)&=-E_y(x,y,-z)\\
E_z(x,y,z)&=\phantom{-}E_z(x,y,-z)\,.
\end{aligned}\right.
\tag{1.46}
\end{equation}
In particolare, quindi, sul piano $z=0$ il campo elettrico (se è definito) ha solo la componente $z$, cioè è ortogonale al piano stesso.

\subsection{Il campo elettrico di distribuzioni simmetriche}

Riprendiamo in considerazione gli esempi del paragrafo precedente e utilizziamo il teorema di Gauss per calcolare il campo elettrico.
\begin{enumerate}
\item Sorgenti a simmetria sferica: $\rho=\rho(r)$, $\vec E(\vec r\,)=E_r(r)\,\hat r$.\\
Presa come \emph{superficie di Gauss} (cioè superficie che sfruttiamo per utilizzare il teorema di Gauss) la sfera di raggio $r$ concentrica con la distribuzione di carica, detta $Q(r)$ la carica contenuta dentro la sfera di raggio $r$, dal teorema di Gauss si ha
\begin{equation}
\Phi(\vec E)=4\pi r^2E_r(r)=4\pi kQ(r)\quad\Rightarrow\quad\vec E(\vec r\,)=k\,\frac{Q(r)}{r^2}\,\hat r
\tag{1.47}
\end{equation}
dove
\begin{equation}
Q(r)\equiv\int_{r'\le r}\rho(r')\,dV=\int_{r'\le r}\rho(r')\,r'^2\,dr'\,d\Omega=4\pi\int_0^r\rho(r')\,r'^2\,dr'
\tag{1.48}
\end{equation}
(l'ultimo passaggio è reso possibile dal fatto che l'integrando $\rho(r')$ non dipende dagli angoli $\theta$ e $\phi$, e l'angolo solido totale $\int d\Omega=4\pi$).\\
Quindi: se $\rho(r)=0$ per $r>R$, il campo elettrico per $r>R$ è identico a quello di una carica puntiforme uguale alla carica totale, posta nel centro di simmetria; se $\rho(r)=0$ per $r<R$, il campo elettrico è nullo per $r<R$ (questo è il caso, per esempio, di una sfera con distribuzione superficiale uniforme $\sigma$ sulla sua superficie); se
\begin{equation}
\rho(r)=\begin{cases}\rho_0\equiv Q/\frac43\pi R^3&r\le R\\[4pt] 0&r>R\end{cases}
\quad\Rightarrow\quad
\vec E(\vec r\,)=\begin{cases}k\,\dfrac{Q}{R^3}\,\vec r&r\le R\\[8pt] k\,\dfrac{Q}{r^2}\,\hat r&r\ge R\end{cases}
\tag{1.49}
\end{equation}
Nelle figure 1.9a e 1.9b è riportato il campo elettrico rispettivamente di una distribuzione uniforme $\sigma$ sulla superficie di una sfera di raggio $R$ e quello di una distribuzione uniforme $\rho_0$ all'interno della stessa sfera.
\begin{center}
\begin{tabular}{cc}
\includegraphics[scale=1.2]{fig1_9a} & \includegraphics[scale=1.2]{fig1_9b}\\
{\small fig.~1.9a} & {\small fig.~1.9b}
\end{tabular}
\end{center}
\item Sorgenti a simmetria cilindrica: $\rho=\rho(\varrho)$ \ $(\varrho\equiv\sqrt{x^2+y^2}\,)$; \ $\vec E(\vec r\,)=E_\varrho(\varrho)\,\hat\varrho$.\\
Presa come superficie di Gauss un cilindro di raggio $\varrho$ e altezza $h$, con asse coincidente con l'asse $z$, il flusso attraverso le basi è nullo, mentre quello attraverso la superficie laterale è $2\pi\varrho h\,E_\varrho$, per cui, se $\lambda(\varrho)$ è la carica per unità di altezza contenuta nel cilindro di raggio $\varrho$ \ $\big(\lambda(\varrho)=Q(\varrho)/h\big)$,
\begin{equation}
\vec E(\vec r\,)=\frac{2k\lambda(\varrho)}{\varrho}\,\hat\varrho\,;\qquad\quad\lambda(\varrho)=2\pi\int_0^\varrho\rho(r)\,r\,dr
\tag{1.50}
\end{equation}
Nel caso particolare di un filo rettilineo con distribuzione lineare costante $\lambda$, $\lambda(\varrho)=\lambda\quad\forall\varrho>0$, e
\begin{equation}
\vec E(\vec r\,)=\frac{2k\lambda}{\varrho}\,\hat\varrho
\tag{1.51}
\end{equation}
e quindi vediamo che il campo diverge come $1/\varrho$. Questo risultato, che abbiamo anticipato nel paragrafo 3, vale per una distribuzione lineare qualsiasi: vedremo nel paragrafo 10 che in ogni punto $P$ di una distribuzione lineare la componente radiale di $\vec E$ diverge come $2k\lambda(P)/\varrho$ (per una distribuzione lineare generica il campo ha -- in generale -- anche altre componenti, ma queste restano finite e continue).
\item Distribuzione con simmetria piana: $\rho=\rho(z)$, \quad $\vec E=\big(0,\,0,\,E_z(z)\big)$.\\
Supponiamo che la distribuzione sia diversa da zero solo fra due piani: $\rho(z)\neq0$ per $z_1\le z\le z_2$. Nella regione priva di sorgenti $z\ge z_2$ (oppure $z\le z_1$) prendiamo come superficie di Gauss un cilindro con le basi ortogonali all'asse $z$ (figura 1.10): il flusso totale è nullo, quello attraverso la superficie laterale anche (tubo di flusso), quindi il flusso entrante in una base ($z=a$) è uguale a quello uscente dall'altra base ($z=b$). Siccome $\vec E$ è ortogonale alle basi e costante su ciascuna di esse, ne segue $E_z(a)=E_z(b)\quad\forall a,\,b\ge z_2$; $\forall a,\,b\le z_1$, cioè il campo è costante sia per $z\ge z_2$ che per $z\le z_1$, in generale con valori diversi nelle due regioni: infatti, se ora prendiamo come superficie di Gauss un cilindro come prima, ma ora con una base nella regione $z\le z_1$ e l'altra nella regione $z\ge z_2$ (figura 1.10), troviamo che
\fig{fig1_10}{1.10}
\begin{equation}
E_z(z\ge z_2)-E_z(z\le z_1)=4\pi k\int_{z_1}^{z_2}\rho(z')\,dz'\equiv4\pi k\sigma
\tag{1.52}
\end{equation}
dove con $\sigma$ abbiamo indicato la carica per unità di area contenuta fra i piani $z=z_1$ e $z=z_2$.\\
In questo caso, siccome $\vec E=\big(0,\,0,\,E_z(z)\big)$, il problema è unidimensionale per cui si può arrivare semplicemente allo stesso risultato anche utilizzando la (1.35) (cioè la forma locale del teorema di Gauss) che assume la forma
\begin{equation}
\dive\vec E=\frac{d\,E_z}{dz}=4\pi k\rho(z).
\tag{1.53}
\end{equation}
Quindi, sappiamo che il campo è costante nelle regioni senza sorgenti, conosciamo la discontinuità quando si passa da una regione all'altra, ma questa volta il teorema di Gauss non ci permette di determinare il campo: tra poco capiremo perché.\\
Supponiamo ora che la nostra distribuzione sia una distribuzione superficiale piana $\sigma$, costante, sul piano $z=0$: in questo caso il gruppo di invarianza delle sorgenti contiene anche la riflessione rispetto al piano $z=0$, come pure le rotazioni di 180$^\circ$ rispetto ad ogni retta giacente in detto piano. Allora l'orbita di ogni punto $P$ è costituita sia dal piano per $P$ \ $z=z_0$ che da quello simmetrico $z=-z_0$, per cui si ottiene che $E_z(-z)=-E_z(z)\quad\forall z\neq0$. Quindi, per la (1.52),
\begin{equation}
\left\{\begin{aligned}
E_z(z)&=\phantom{-}2\pi k\sigma\quad\text{per }z>0\\
E_z(z)&=-2\pi k\sigma\quad\text{per }z<0
\end{aligned}\right.
\tag{1.54}
\end{equation}
Quindi la distribuzione superficiale produce una discontinuità pari a $4\pi k\sigma$ nella
componente normale del campo (in questo caso l'unica): anche questo, come già anticipato, è un risultato generale valido per qualsiasi distribuzione superficiale. Per esempio, se abbiamo una distribuzione costante $\sigma$ sulla superficie di una sfera di raggio $R$, sappiamo \big(dalla (1.47)\big) che dentro la sfera il campo è nullo, mentre all'esterno è
\begin{equation}
\vec E(\vec r\,)=k\,\frac{4\pi\sigma R^2}{r^2}\,\hat r\quad r>R
\tag{1.55}
\end{equation}
e vediamo, anche in questo caso, che la componente radiale del campo subisce una discontinuità pari a $4\pi k\sigma$ attraverso la distribuzione superficiale.\\
Le linee di forza escono dal piano (da entrambe le parti) se $\sigma>0$, entrano se $\sigma<0$.\\
Supponiamo ora di avere due piani paralleli rispettivamente con distribuzioni costanti $\sigma$ e $-\sigma$. Il campo elettrico si ottiene sommando i campi prodotti da ciascuna delle due distribuzioni (principio di sovrapposizione): i due campi sono concordi fra i due piani e discordi nelle due regioni esterne (nella figura 1.11 le frecce superiori rappresentano il campo generato dalla distribuzione positiva, quelle inferiori quello generato dalla distribuzione negativa), quindi
\fig{fig1_11}{1.11}
\begin{equation}
\left\{\begin{aligned}
&E=4\pi k\sigma&&\text{fra i due piani}\\
&E=0&&\text{all'esterno}
\end{aligned}\right.
\tag{1.56}
\end{equation}
Un sistema del genere è detto \emph{condensatore}, in quanto ``condensa'' il campo nella regione compresa fra i due piani (due piani paralleli con distribuzioni uguali possiamo quindi chiamarlo un ``anticondensatore'', in quanto il campo è diverso da zero solo all'esterno).\\
Torniamo alla distribuzione $\rho(z)$ fra i due piani $z=z_1$ e $z=z_2$. Possiamo dividere l'intervallo $z_1$, $z_2$ in $n$ strati di ampiezza $\Delta z$, e considerare il campo elettrico come la sovrapposizione dei campi prodotti da questi $n$ strati: per $n\to\infty$ ciascuno di questi strati finisce per generare il campo di una distribuzione superficiale piana di densità $\rho(z_i)\Delta z$, quindi campi opposti per $z>z_2$ e per $z<z_1$. Di conseguenza anche la sovrapposizione di tutti questi campi ha la stessa proprietà: $E_z(z\ge z_2)=-E_z(z\le z_1)$. Questa è l'informazione che ci mancava per determinare il campo della distribuzione $\rho(z)$: se non ci sono altre sorgenti
\begin{equation}
E_z(z)=\begin{cases}
\phantom{-}2\pi k\displaystyle\int_{z_1}^{z_2}\rho(z')\,dz'&\text{per }z\ge z_2\\[12pt]
-2\pi k\displaystyle\int_{z_1}^{z_2}\rho(z')\,dz'&\text{per }z\le z_1
\end{cases}
\tag{1.57}
\end{equation}
Abbiamo detto ``se non ci sono altre sorgenti''; questa è la ragione per cui simmetrie e teorema di Gauss non sono stati sufficienti a determinare il campo: simmetrie e teorema di Gauss sono compatibili con la presenza di altre sorgenti lontane con le stesse simmetrie, per esempio la distribuzione $\rho(z)$ potrebbe trovarsi all'interno di un condensatore. In effetti la (1.52) determina il campo a meno di un campo costante parallelo all'asse $z$.
\end{enumerate}
Un'ultima osservazione: può venire il sospetto che la ragione per cui il campo elettrico di una distribuzione piana uniforme $\sigma$ è costante dalle due parti della distribuzione possa scaturire da una qualche proprietà di invarianza delle sorgenti. Dovrebbe trattarsi di una proprietà peculiare della legge $1/r^2$, perché noi la abbiamo dedotta usando il teorema di Gauss: la risposta è sì, ed ha le sue radici nel fatto che il campo di una carica puntiforme resta invariato se $q\to\lambda^2q$ e $\vec r\to\lambda\vec r$, \ $\forall\lambda>0$: una carica di 1\,C produce alla distanza di 1\,m lo stesso campo che una carica di 100\,C produce alla distanza di 10\,m. Questa è quella che si chiama \emph{invarianza di scala}: se tutte le lunghezze vengono cambiate per un fattore e le cariche per lo stesso fattore al quadrato, il campo elettrico resta invariato. Ma in questo modo una distribuzione uniforme su un piano resta inalterata ($\sigma=q/\Delta S\equiv\lambda^2q/\lambda^2\Delta S$), e di conseguenza $\vec E(\vec r\,)=\vec E(\lambda\vec r\,)$, cioè genera un campo indipendente dalla distanza (se invece la superficie non è piana le sue dimensioni cambiano).

\subsection{Il potenziale del campo elettrico}

Il campo elettrico di una carica puntiforme è un campo centrale, quindi conservativo; grazie al principio di sovrapposizione è conservativo anche il campo elettrico di più cariche puntiformi, e quello delle distribuzioni continue (le cariche devono essere in quiete). Il potenziale del campo di una carica puntiforme è dato da (lo si confronti con il potenziale gravitazionale dovuto a una massa puntiforme)
\begin{equation}
\varphi(r)=k\,\frac{q}{r}
\tag{1.58}
\end{equation}
avendo posto $\varphi(\infty)=0$. Per una distribuzione di volume
\begin{equation}
\varphi(\vec r\,)=k\int_V\frac{\rho(\vec r\,')}{|\vec r-\vec r\,'|}\,dV
\tag{1.59}
\end{equation}
e analoghe per le distribuzioni superficiali e per quelle lineari. Si noti che siccome il campo elettrico è un campo di forza per unità di carica, anche il potenziale elettrico è l'energia potenziale per unità di carica.\\
Richiamiamo le proprietà fondamentali dei campi conservativi, con riferimento al campo elettrico.
\begin{itemize}
\item[i)] Dato un qualsiasi circuito chiuso $\gamma$ orientato, si definisce \textbf{circuitazione} del campo elettrico (la definizione si applica a qualsiasi campo vettoriale, anche non conservativo)
\begin{equation}
\oint_\gamma\vec E\cdot d\vec s\defeq\int_\gamma E_t(s)\,ds
\tag{1.60}
\end{equation}
dove $E_t(s)$ è, nel punto di ascissa curvilinea $s$ di $\gamma$, la componente di $\vec E$ tangente alla curva: $E_t\equiv\vec E\cdot\hat\tau$ ($\hat\tau$ è il versore tangente alla curva) e abbiamo introdotto la definizione di elemento di linea orientata $d\vec s\equiv\hat\tau\,ds$. Siccome il campo elettrico è conservativo la sua circuitazione è sempre nulla
\begin{equation}
\oint_\gamma\vec E\cdot d\vec s=0
\tag{1.61}
\end{equation}
e l'integrale di linea fra due punti $P_1$ e $P_2$ è indipendente dal percorso.
\item[ii)] Il potenziale è definito a meno di una costante additiva: se $P_0$ è un punto arbitrario
\begin{equation}
\varphi(P)=-\int_{P_0}^P\vec E\cdot d\vec s.
\tag{1.62}
\end{equation}
Quindi, se le sorgenti sono fisse e $q$ è una carica ``esploratrice'', $-q\varphi(P)$ è il lavoro fatto dalle forze del campo elettrico sulla carica $q$ quando questa viene spostata dal punto $P_0$ al punto $P$: $q\varphi(P)$ è quindi il lavoro (positivo o negativo) che occorre fare \underline{dall'esterno} (cioè \underline{contro} le forze del campo) per portare la carica da $P_0$ a $P$. In generale $q\varphi(P_2)-q\varphi(P_1)$ è il lavoro che occorre fare per portare la carica dal punto $P_1$ al punto $P_2$.
\item[iii)] La relazione fra il campo elettrico e il potenziale è uguale a quella che si ha fra forza e (energia) potenziale:
\begin{equation}
\vec E(x,y,z)=-\grad\varphi(x,y,z)\equiv-\Big(\frac{\partial\varphi}{\partial x},\,\frac{\partial\varphi}{\partial y},\,\frac{\partial\varphi}{\partial z}\Big)\,\cdot
\tag{1.63}
\end{equation}
L'operatore \emph{gradiente} (grad) sovente viene anche indicato con il simbolo $\nabla$ (``nabla'')
\begin{equation}
\grad\equiv\nabla\equiv\Big(\frac{\partial}{\partial x}\,,\,\frac{\partial}{\partial y}\,,\,\frac{\partial}{\partial z}\Big)\,\cdot
\tag{1.64}
\end{equation}
Poiché (teorema di Schwarz) per ogni funzione (con derivate seconde ``miste'' continue) $\dfrac{\partial^2f}{\partial x\,\partial y}=\dfrac{\partial^2f}{\partial y\,\partial x}$\,, definita la \textbf{rotazione} (o \textbf{rotore}) del campo elettrico (o di qualunque altro campo vettoriale) come
\begin{equation}
\rot\vec E\defeq\Big(\frac{\partial E_z}{\partial y}-\frac{\partial E_y}{\partial z}\,,\,\frac{\partial E_x}{\partial z}-\frac{\partial E_z}{\partial x}\,,\,\frac{\partial E_y}{\partial x}-\frac{\partial E_x}{\partial y}\Big)\equiv\nabla\wedge\vec E
\tag{1.65}
\end{equation}
dalla (1.63) si ha
\begin{equation}
\rot\vec E(x,y,z)=0
\tag{1.66}
\end{equation}
per cui si dice che il campo elettrico è un campo \emph{irrotazionale}. Nel capitolo 7 vedremo il significato geometrico della rotazione di un campo vettoriale; la parola stessa fa capire che ha a che fare, nel caso in cui il campo vettoriale sia il campo delle velocità di un fluido, con i moti rotatori del fluido: un campo di velocità è irrotazionale se non ci sono vortici.\\
La (1.66) è condizione necessaria perché un campo vettoriale sia conservativo; è anche sufficiente se lo spazio è semplicemente connesso (cioè ogni circuito chiuso è riducibile con continuità ad un punto): quando discuteremo i campi variabili (nel tempo) incontreremo casi in cui la (1.66) è soddisfatta, ma il campo elettrico non è conservativo.
\item[iv)] Il campo elettrico è in ogni punto (non singolare) ortogonale alla superficie equipotenziale passante per quel punto ed è diretto nel verso in cui il potenziale decresce. Le superfici equipotenziali integrano le informazioni grafiche sul campo elettrico fornite dalle linee di forza: siccome le superfici equipotenziali di una carica puntiforme sono sfere, anche in presenza di più sorgenti -- analogamente a quanto già
visto per le linee di forza -- le superfici equipotenziali tendono alla forma sferica al diminuire della distanza da una carica puntiforme.
\end{itemize}
Nel sistema di Gauss il potenziale si misura in \texttt{\textbf{statvolt}}: fra due punti si ha una differenza di potenziale di 1\,statvolt se il lavoro che occorre fare per portare una carica di 1\,franklin da un punto all'altro è di 1\,erg. Nel SI l'unità di potenziale è il più noto volt (\texttt{V}): 1\,volt $=$ 1\,joule/1\,coulomb. Si ha
\begin{equation}
1\,\mathrm{V}=\frac{10^7\,\mathrm{erg}}{3\times10^9\,\mathrm{franklin}}=\frac{1}{300}\,\mathrm{statvolt}\,;\qquad\qquad1\,\mathrm{statvolt}=300\,\mathrm{V}
\tag{1.67}
\end{equation}
Siccome i potenziali (o meglio: le differenze di potenziale) sono misurabili molto più direttamente e semplicemente che non le forze e le cariche, ed inoltre i campi elettrici di solito vengono generati stabilendo una differenza di potenziale fra due conduttori ad una certa distanza, si preferisce esprimere la misura dei campi elettrici in volt/metro o in statvolt/cm (anziché in newton/coulomb o dyne/franklin): naturalmente si tratta della stessa unità di misura espressa in termini di grandezze diverse (talvolta si usa il V/cm che è 100 volte il V/m e 1/300\,ues).\\
In fisica atomica viene usata normalmente come unità di energia l'elettronvolt (\texttt{\textbf{eV}}): come dice il nome è la carica dell'elettrone per 1 volt:
\begin{equation}
\begin{split}
1\,\mathrm{eV}&=1.6\times10^{-19}\,\mathrm{coulomb}\cdot1\,\mathrm{volt}=1.6\times10^{-19}\,\mathrm{joule}=\\
&=1.6\times10^{-12}\,\mathrm{erg}\,.
\end{split}
\tag{1.68}
\end{equation}

\paragraph{Esempi}
\begin{enumerate}
\item Consideriamo un elettrone che gira attorno al protone a distanza di $0.53\times10^{-8}$\,cm (modello di Rutherford dell'atomo di idrogeno). Il campo elettrico a cui è soggetto l'elettrone è
\[
E=\frac{e}{r^2}=\frac{4.8\times10^{-10}}{\big(0.53\times10^{-8}\big)^2}=1.7\times10^7\,\mathrm{ues}=5.1\times10^9\,\mathrm{V/cm}.
\]
Il potenziale (rispetto all'infinito) nella posizione dell'elettrone è
\[
\varphi=\frac{e}{r}=\frac{4.8\times10^{-10}}{0.53\times10^{-8}}=0.09\,\mathrm{ues}=27.2\,\mathrm{V}
\]
e quindi l'energia potenziale dell'elettrone è $U=-e\varphi=-27.2\,\mathrm{eV}$. Siccome per le orbite circolari nei campi di forza proporzionali a $1/r^2$ l'energia cinetica è (in valore assoluto) metà dell'energia potenziale, l'energia totale dell'elettrone è di $-13.6\,\mathrm{eV}$: $13.6\,\mathrm{eV}$ è l'energia necessaria per ionizzare l'atomo di idrogeno (le energie di 1$^{\rm a}$ ionizzazione degli atomi, cioè l'energia necessaria per estrarre un solo elettrone, vanno dai (circa) 4\,eV del Cs ai 25\,eV dell'He).
\item Potenziale dovuto ad una sfera di raggio $R$ con carica $Q$ distribuita a simmetria sferica. Il campo elettrico è dato dalla (1.47) e all'esterno dipende solo dalla carica totale. Preso $\varphi(\infty)=0$ (tutte le volte che ciò è possibile faremo tacitamente questa scelta), si ha
\begin{equation}
\varphi(r)=k\,\frac{Q}{r}\qquad\quad r\ge R\,.
\tag{1.69}
\end{equation}

Se la carica è (uniformemente) distribuita sulla superficie della sfera il campo elettrico all'interno è nullo e quindi il potenziale è costante: anche se il campo elettrico è discontinuo (v.\ figura 1.9a) attraverso la superficie della sfera, il potenziale (ma non le sue derivate) è continuo (se l'integrando è discontinuo ma finito, l'integrale è continuo), e quindi all'interno vale $kQ/R$ (figura 1.12a).
\begin{center}
\begin{tabular}{cc}
\includegraphics[scale=1.2]{fig1_12a} & \includegraphics[scale=1.2]{fig1_12b}\\
{\small fig.~1.12a} & {\small fig.~1.12b}
\end{tabular}
\end{center}
Se invece la carica è uniformemente distribuita all'interno della sfera, il campo elettrico è dato dalla (1.49) ed il potenziale (all'interno) per la (1.62) è
\begin{equation}
\begin{split}
\varphi(r)&=\varphi(R)-k\int_R^r\frac{Q}{R^3}\,\vec r\,'\cdot d\vec r\,'=k\,\frac{Q}{R}-k\int_R^r\frac{Q}{R^3}\,r'\,dr'=\\
&=k\,\frac{Q}{R}-k\,\frac{Q}{2R^3}(r^2-R^2)=k\,\frac{Q}{2R}\Big(3-\frac{r^2}{R^2}\Big)\qquad r\le R\,.
\end{split}
\tag{1.70}
\end{equation}
Nella figura 1.12b è riportato $\varphi(r)$, nel caso di una distribuzione $\rho>0$ (altrimenti il grafico deve essere ribaltato).
\item Condensatore piano: distribuzioni piane uniformi $\sigma$ e $-\sigma$ a distanza $d$. Il campo elettrico è dato dalla (1.56): siccome all'esterno il campo è nullo, il potenziale nelle due regioni esterne è costante, con valori diversi. Preso il piano $z=0$ coincidente con la distribuzione $-\sigma$ e $z=d$ quello della distribuzione $\sigma$, si ha (figura 1.13)
\fig{fig1_13}{1.13}
\begin{equation}
\varphi(z)=\begin{cases}0&\text{per }z\le0\\ 4\pi k\sigma\,z&\text{per }0\le z\le d\\ 4\pi k\sigma\,d&\text{per }z\ge d\end{cases}
\tag{1.71}
\end{equation}
dove la scelta $\varphi=0$ per $z\le0$ è arbitraria: la cosa importante è che la differenza di potenziale fra le due distribuzioni è $4\pi k\sigma d\equiv Ed$.
\item Filo rettilineo indefinito con distribuzione lineare costante $\lambda$. Il campo elettrico è dato dalla (1.51). Preso un punto a distanza $\varrho_0$ dal filo, si ha
\begin{equation}
\varphi(\varrho)=-\int_{\varrho_0}^\varrho\frac{2\lambda k}{\varrho'}\,d\varrho'=-2k\lambda\,\log\frac{\varrho}{\varrho_0}\,\cdot
\tag{1.72}
\end{equation}
La scelta di $\varrho_0$, cioè del punto in cui il potenziale si annulla, è arbitraria (purché al finito e $\neq0$), in quanto influisce solo come una costante additiva al potenziale. Come si vede, in questo caso non è possibile porre $\varphi(\infty)=0$: se $\varrho_0\to\infty$ il potenziale diverge. Questo perché il campo elettrico va come $1/\varrho$ a grandi distanze e ciò è dovuto al fatto che la distribuzione si estende su una regione infinita. In questo caso, l'espressione analoga alla (1.59)
\[
\varphi(\vec r\,)=k\int\frac{\lambda}{|\vec r-\vec r\,'|}\,dl
\]
non è corretta in quanto presuppone $\varphi(\infty)=0$, e in effetti l'integrale diverge: il modo corretto di scrivere la (1.59) è
\begin{equation}
\varphi(\vec r\,)-\varphi(\vec r_0)=k\int\rho(\vec r\,')\Big(\frac{1}{|\vec r-\vec r\,'|}-\frac{1}{|\vec r_0-\vec r\,'|}\Big)\,dV
\tag{1.73}
\end{equation}
e espressioni analoghe per distribuzioni superficiali e lineari. Se la distribuzione è contenuta in una regione finita (non è però una condizione necessaria) l'integrale della differenza è la differenza degli integrali e ciascuno di essi tende a zero quando $r$ (o $r_0$) tende all'infinito.
\end{enumerate}

\subsection{Campo elettrico e potenziale a grandi distanze}

Supponiamo che le sorgenti siano contenute in un volume (finito) $V$, per cui possiamo utilizzare la (1.59) per determinare il comportamento del potenziale (e quindi del campo elettrico) a grandi distanze. Con le notazioni della (1.59) ($\vec r$ individua il punto ``potenziato'', $\vec r\,'$ il generico punto nella regione in cui sono presenti le sorgenti, detto anche punto ``potenziante''), si ha
\begin{equation}
|\vec r-\vec r\,'|^{-1}=(r^2-2\vec r\cdot\vec r\,'+r'^2)^{-1/2}=\frac1r\Big(1-2\,\frac{\vec r\cdot\vec r\,'}{r^2}+\frac{r'^2}{r^2}\Big)^{-1/2}.
\tag{1.74}
\end{equation}
Presa l'origine delle coordinate in un punto prossimo alla distribuzione delle sorgenti, se $r\gg r'\quad\forall\vec r\,'\in V$, siccome per $|\delta|<1$, \ $(1+\delta)^{-1/2}=1-\frac12\delta+\frac38\delta^2+O(\delta^3)$, posto
\[
\delta=-2\,\frac{\vec r\cdot\vec r\,'}{r^2}+\frac{r'^2}{r^2}\equiv-2\,\frac{r'}{r}\cos\theta'+\frac{r'^2}{r^2}
\]
($\theta'$ è l'angolo fra $\vec r$ e $\vec r\,'$), dalla (1.74), tenendo i termini solo fino all'ordine $(r'/r)^2$, si ha
\begin{equation}
|\vec r-\vec r\,'|^{-1}=\frac1r\Big(1+\frac{r'}{r}\cos\theta'+\frac{r'^2}{r^2}\,\frac{3\cos\theta'^2-1}{2}+O\big(\frac{r'^3}{r^3}\big)\Big)
\tag{1.75}
\end{equation}
La (1.75), inserita nella (1.59), fornisce lo sviluppo del potenziale come serie di potenze in $1/r$, valido per $r\gg r'\quad\forall\vec r\,'\in V$:
\begin{equation}
\varphi(\vec r\,)=k\,\frac1r\int_V\rho(\vec r\,')\,dV+k\,\frac{1}{r^2}\int_Vr'\cos\theta'\rho(\vec r\,')\,dV+\cdots
\tag{1.76}
\end{equation}
Il primo integrale è la carica totale $q$; il 2$^\circ$ termine, essendo $\cos\theta'\equiv\dfrac{\vec r\cdot\vec r\,'}{rr'}$ e definito il \textbf{momento di dipolo} $\vec p$ della distribuzione come
\begin{equation}
\vec p\defeq\int_V\vec r\,'\rho(\vec r\,')\,dV
\tag{1.77}
\end{equation}
può essere scritto $k\,\dfrac{\vec p\cdot\vec r}{r^3}$. In definitiva
\begin{equation}
\varphi(\vec r\,)=k\Big[\,\frac qr+\frac{\vec p\cdot\vec r}{r^3}+\frac12\sum_{ij}\frac{Q_{ij}x_ix_j}{r^5}+\cdots\Big]
\tag{1.78}
\end{equation}
dove per completezza (o per curiosità) abbiamo riportato anche l'espressione del 3$^\circ$ termine (termine di \textbf{quadrupolo}), con
\begin{equation}
Q_{ij}\defeq\int_V(3x'_ix'_j-r'^2\delta_{ij})\rho(\vec r\,')\,dV
\tag{1.79}
\end{equation}
(espressioni analoghe, anche se più complesse, per i termini successivi).\\
La (1.78) è nota come lo sviluppo del potenziale in \emph{termini di multipolo}: termine di monopolo (cioè la carica), di dipolo, quadrupolo, ottupolo eccetera.\\
A grandi distanze il comportamento del potenziale è determinato dal primo termine non nullo: se $q\neq0$, il potenziale va come $1/r$, quindi il campo elettrico come $1/r^2$; se la carica totale è nulla, ma non il momento di dipolo, il potenziale va come $1/r^2$ ed inoltre dipende dalla direzione:
\[
k\,\frac{\vec p\cdot\vec r}{r^3}=k\,\frac{p\cos\theta}{r^2}
\]
ed il campo elettrico, detto \textbf{campo di dipolo}, va come $1/r^3$, e così via per i multipoli superiori.\\
La (1.77) definisce il momento di dipolo per una distribuzione di volume: espressioni analoghe per distribuzioni superficiali o lineari. Nel caso di cariche puntiformi si ha
\begin{equation}
\vec p=\sum_i^nq_i\vec r_i.
\tag{1.80}
\end{equation}
Il momento di dipolo (e anche i multipoli superiori) di una distribuzione di cariche (continua e/o discreta), in generale dipende dalla scelta del punto, detto \textbf{polo}, rispetto al quale i vettori $\vec r_i$ nella (1.80) \big(o il vettore $\vec r\,'$ nella (1.77)\big) individuano la posizione delle cariche: se nella (1.80) $\vec r\,'=\vec r-\vec a$ (cioè il polo viene spostato dall'origine al punto $\vec r=\vec a$), si ha
\begin{equation}
\vec p\,'=\sum_i^nq_i\vec r\,'_i=\vec p-q\,\vec a
\tag{1.81}
\end{equation}
La (1.81) non solo dice come dipende dalla scelta del polo il momento di dipolo, ma soprattutto dice che se la carica totale del sistema è nulla, il momento di dipolo del sistema è indipendente dal polo, e quindi è una proprietà intrinseca del sistema. Lo stesso risultato è vero per i multipoli superiori: il primo multipolo non nullo è indipendente dal polo.\\
Il prototipo di sistema con momento (intrinseco) di dipolo è costituito da due cariche $q$ e $-q$. Se $\vec d$ è la posizione della carica positiva rispetto a quella negativa, il momento di dipolo del sistema è dato da
\fig{fig1_14}{1.14}
\begin{equation}
\vec p=q\,\vec d
\tag{1.82}
\end{equation}
ed il potenziale a grandi distanze ($r\gg d$) è dato dal 2$^\circ$ termine della (1.78).\\
Si definisce \textbf{dipolo puntiforme} un sistema per il quale il potenziale è dato dal 2$^\circ$ termine della (1.78) a qualunque distanza: esso si ottiene, per esempio, dal dipolo fisico costituito da due cariche a distanza $d$ prendendo il limite $d\to0$, $q\to\infty$ con la condizione $q\,\vec d\equiv\vec p$ (costante); nel limite $d\to0$ qualunque distanza è grande rispetto a $d$. Il dipolo puntiforme è ovviamente solo una schematizzazione matematica (come la carica puntiforme).

Il prototipo di quadrupolo è costituito (come dice il nome) da quattro cariche: deve essere nulla la carica totale e anche il momento di dipolo; esso è quindi costituito da due dipoli opposti in posizioni diverse. Ci sono diversi modi di disporli: per esempio allineati (come nella parte sinistra della figura 1.15), oppure affiancati (come nella parte destra); ciò è dovuto alla maggior complessità del momento di quadrupolo (un tensore del 2$^\circ$ ordine) rispetto al dipolo (un vettore).
\fig{fig1_15}{1.15}
E così via: per esempio (figura 1.16), quattro cariche $+q$ e quattro $-q$ disposte alternativamente ai vertici di un cubo costituiscono un ottupolo (due facce opposte del cubo costituiscono due quadrupoli opposti), ed il sistema dà luogo ad un potenziale che a grandi distanze va come $1/r^4$, con una complicata dipendenza dalla direzione.
\fig{fig1_16}{1.16}
Infine, se $q\neq0$ e si vuole ottenere il potenziale (a grandi distanze) fino all'ordine $1/r^2$ compreso, è possibile sfruttare la dipendenza dal polo del momento di dipolo per far sì che il 2$^\circ$ termine nella (1.78) sia nullo, per cui il primo termine da solo fornisce il potenziale correttamente a meno di termini $1/r^3$: grazie alla (1.81), se si prende l'origine nel punto $\vec r=\vec p/q$, il momento di dipolo si annulla. Così, per esempio, nel caso di due cariche uguali, se l'origine viene presa nel punto di mezzo fra le due cariche, $k\,q/r$ differisce dal potenziale $\varphi(\vec r\,)$ solo per termini dell'ordine $1/r^3$.

\subsection{Dal campo alle sorgenti: le discontinuità}

Date le sorgenti, siano esse cariche puntiformi, distribuzioni lineari, superficiali o di volume, il campo elettrico ed il potenziale sono ben determinati. Viceversa, se è dato il campo elettrico, o il potenziale, è possibile determinare tutte le sorgenti del campo.\\
La distribuzione di volume si determina semplicemente calcolando la divergenza di $\vec E$ oppure, se è dato il potenziale, per la (1.63) e (1.35)
\begin{equation}
\dive\big(\grad\varphi(x,y,z)\big)\equiv\frac{\partial^2\varphi}{\partial x^2}+\frac{\partial^2\varphi}{\partial y^2}+\frac{\partial^2\varphi}{\partial z^2}=-4\pi k\,\rho(x,y,z)
\tag{1.83}
\end{equation}
cioè, definito l'\textbf{operatore di Laplace} (o \textbf{laplaciano})
\begin{equation}
\Delta\defeq\Big(\frac{\partial^2}{\partial x^2}+\frac{\partial^2}{\partial y^2}+\frac{\partial^2}{\partial z^2}\Big)
\tag{1.84}
\end{equation}
calcolando il laplaciano del potenziale:
\begin{equation}
\Delta\,\varphi(x,y,z)=-4\pi k\,\rho(x,y,z)
\tag{1.85}
\end{equation}
La (1.85) è nota come \textbf{equazione di Poisson} ed è l'equazione fondamentale dell'elettrostatica, in quanto \textit{i}) $\varphi$ è il potenziale del campo elettrico, quindi tiene conto del fatto che il campo elettrico è conservativo; \textit{ii}) è equivalente alla (1.35), cioè alla 1$^{\rm a}$ equazione di Maxwell.\\
La (1.85) è valida in tutti i punti in cui il campo non è singolare o presenti discontinuità: le eventuali cariche puntiformi sono localizzate nei punti $\vec r_i$ in cui il campo elettrico diverge come $1/|\vec r-\vec r_i|^2$ (oppure il potenziale diverge come $1/|\vec r-\vec r_i|$), e
per determinare il valore di $q_i$ basta calcolare il flusso del campo elettrico uscente da una sferetta centrata in $\vec r_i$, nel limite di raggio che tende a zero (il limite è necessario per escludere altre cariche, discrete o distribuite), oppure moltiplicare il potenziale per $|\vec r-\vec r_i|$ e prendere il limite per $\vec r\to\vec r_i$.

\paragraph{Esempio}
\begin{quote}
È dato il potenziale $\varphi(r)=q\,\dfrac{e^{-\alpha r}}{r}\Big(1+\dfrac{\alpha r}{2}\Big)$. Vogliamo determinare le sorgenti.\\
C'è una singolarità $1/r$ nell'origine e siccome per $r\to0$ \ $\lim r\varphi(r)=q$, \ la carica è $q$. Per determinare la distribuzione di volume calcoliamo la divergenza di $\vec E$: siccome il potenziale dipende solo da $r$ il campo elettrico è radiale e la sua divergenza (ovvero il laplaciano di $\varphi$) è una funzione solo di $r$; in casi come questi (cioè per campi radiali) dalla (1.24) (teorema della divergenza), preso come volume la sfera di raggio $r$, si ha
\[
4\pi r^2E_r(r)=\int_V\dive\vec E\,dV=4\pi\int_0^r\dive\vec E\,r'^2\,dr'
\]
e presa la derivata di ambo i membri rispetto ad $r$ si ottiene
\begin{equation}
\dive\vec E=\frac{1}{r^2}\,\frac{d}{dr}(\,r^2E_r)\qquad\qquad\text{(per un campo radiale)}
\tag{1.86}
\end{equation}
Siccome
\[
E_r=-\varphi'(r)=q\,e^{-\alpha r}\Big(\frac{1}{r^2}+\frac\alpha r+\frac{\alpha^2}{2}\Big)\quad\Rightarrow\quad\rho(r)=-\frac{q\alpha^3}{8\pi}\,e^{-\alpha r}\quad(r>0).
\]
Si noti che siccome il campo elettrico decresce esponenzialmente all'infinito, la carica totale è nulla (come si può anche verificare direttamente), in quanto il flusso uscente da una sfera di raggio $R$ tende a zero quando $R\to\infty$: la superficie cresce come $R^2$, ma il campo decresce più rapidamente di $R^{-2}$.
\end{quote}
Le distribuzioni superficiali si determinano dal fatto che producono la più volte citata discontinuità nella componente del campo elettrico normale alla superficie:
\begin{equation}
\Delta E_n(P)=4\pi k\sigma(P)
\tag{1.87}
\end{equation}
che si dimostra nel modo seguente: prendiamo come superficie di Gauss un cilindretto contenente il punto $P$ ed avente le basi parallele al piano tangente alla distribuzione superficiale nel punto $P$ (figura 1.17); le basi abbiano raggio $\epsilon a$ e l'altezza del cilindretto sia $\epsilon^2a$, in modo che (quando $\epsilon\to0$) il flusso attraverso la superficie laterale sia infinitesimo rispetto a quello attraverso le basi. Se $\hat n$ è la normale alla distribuzione superficiale nel punto $P$, orientata a piacere, detta $\Delta S$ la superficie delle basi, $\vec E^+$ e $\vec E^-$ rispettivamente il campo sulla base dalla parte di $\hat n$ e su quella dalla parte opposta, dal teorema di Gauss si ha $E_n^+\Delta S-E_n^-\Delta S\equiv\Delta E_n\,\Delta S=4\pi k\sigma(P)\Delta S$ da cui la (1.87). Chiaramente la (1.87) è indipendente dall'orientamento di $\hat n$ perché se $\hat n\to-\hat n$, \ $E_n^+\leftrightarrow-E_n^-$.
\fig{fig1_17}{1.17}
Le componenti di $\vec E$ tangenti alla superficie $S$ sono continue: si consideri un circuito rettangolare avente due lati di lunghezza $\Delta l=\epsilon a$ paralleli al piano tangente ad $S$ in $P$ e da parti opposte, e gli altri due lati di lunghezza $\epsilon^2a$; la circuitazione di $\vec E$ è nulla perché il campo elettrico è conservativo, i lati corti (al limite) non contribuiscono, e quindi $E_t^+\Delta l-E_t^-\Delta l=0$, da cui la tesi.
\fig{fig1_18}{1.18}
Anche per le distribuzioni lineari abbiamo già citato (e incontrato) il comportamento singolare del campo, e quindi anche del potenziale \big(vedi la (1.72)\big). Per vedere che la singolarità $2\lambda(P)/\varrho$ della componente radiale del campo è indipendente dalla forma e dalla lunghezza della curva supporto della distribuzione possiamo ragionare nel modo seguente: in primo luogo ci domandiamo qual è la causa della singolarità $1/\varrho$ nella (1.51), perché questa singolarità potrebbe essere dovuta o al fatto che la distribuzione è lineare (cioè concentrata), oppure al fatto che si estende all'infinito (in generale gli integrali -- in questo caso quello che compare nella (1.17) -- possono divergere o perché diverge l'integrando oppure perché si estendono all'infinito); escludiamo la seconda possibilità perché dalla (1.50) si vede che se la distribuzione non è concentrata, $\vec E$ non diverge. Allora, presa una distribuzione lineare qualsiasi, consideriamo un punto $P$ su di essa; se rimuoviamo un tratto $\Delta l$ della distribuzione attorno al punto $P$, il campo elettrico in $P$ è finito perché le sorgenti sono a distanza non nulla da $P$: quindi la singolarità è dovuta al tratto $\Delta l$, comunque piccolo esso sia. Allora possiamo sostituire alla parte della distribuzione esterna a $\Delta l$ un filo rettilineo con distribuzione uguale a $\lambda(P)$ e avente come direzione quella del tratto (infinitesimo) $\Delta l$. Quindi nel punto $P$ la componente radiale del campo diverge come $2\lambda(P)/\varrho$.

\paragraph{Esempi}
\begin{enumerate}
\item Il potenziale è \ $\varphi=Ae^{-\alpha|x|}\cos\alpha y\quad\alpha>0$. \ Determiniamo le sorgenti.
\[
E_x=\begin{cases}\phantom{-}\alpha Ae^{-\alpha x}\cos\alpha y&\text{per }x>0\\ -\alpha Ae^{\alpha x}\cos\alpha y&\text{per }x<0\end{cases}
\qquad E_y=\alpha Ae^{-\alpha|x|}\sen\alpha y\,;\quad E_z=0.
\]
Per $x\neq0$ \ $\dive\vec E=0$ \ quindi le sorgenti sono nel piano $x=0$.
\[
4\pi k\sigma=E_x(0^+)-E_x(0^-)=2\alpha A\cos\alpha y\,;\qquad\sigma=\frac{\alpha A}{2\pi k}\cos\alpha y.
\]
\item $\varphi(x,y,z)=A\log\dfrac{z+c+\sqrt{x^2+y^2+(z+c)^2}}{z-c+\sqrt{x^2+y^2+(z-c)^2}}$ \quad $(c>0)$ \quad è il potenziale di una distribuzione lineare $\lambda$ sul segmento $x=y=0$, $|z|\le c$. Vogliamo calcolare $\lambda$.\\
In prossimità di un punto $P$ di una distribuzione lineare $\lambda$ di carica, il potenziale è $\varphi=-2\lambda(P)\log\varrho+{}$ termini non divergenti per $\varrho\to0$, dove $\varrho$ è la distanza da $P$ sul piano ortogonale alla distribuzione in $P$. Quindi occorre studiare il comportamento di $\varphi(x,y,z)$ per $|z|<c$ e $\varrho\to0$:
\[
\begin{split}
\varphi&=A\log\frac{z+c+\sqrt{\varrho^2+(z+c)^2}}{z-c+\sqrt{\varrho^2+(z-c)^2}}\approx A\log\frac{2(z+c)}{z-c-(z-c)\big(1+\frac12\varrho^2/(z-c)^2\big)}\\
&=-2A\log\varrho+\text{termini finiti}\quad\Rightarrow\quad\lambda(z)=A\,.
\end{split}
\]
\end{enumerate}

\subsection{L'energia di un sistema di cariche}

Non solo il campo di forze che un sistema di cariche fisse esercita su una carica ``esploratrice'' è conservativo, ma siccome sono conservative le forze (coulombiane) fra ogni coppia di cariche di un sistema, un insieme di cariche è un sistema conservativo. Ciò significa che il lavoro per portare un sistema di cariche da una configurazione ad un'altra è indipendente dalle traiettorie delle singole cariche.\\
Quindi si può definire l'\emph{energia $U$ di un sistema di cariche} in una certa configurazione (o \emph{energia elettrostatica}) come il lavoro che occorre fare per portare le cariche nella data configurazione, a partire dalla configurazione in cui le cariche sono all'infinito e infinitamente lontane fra di loro.\\
Cominciamo con un sistema di $n$ cariche puntiformi. Per portare la 1$^{\rm a}$ carica nella sua posizione finale non si fa nessun lavoro, mentre per portare la 2$^{\rm a}$ si fa il lavoro $k\,q_1q_2/r_{12}$, dove $r_{12}$ è la distanza fra le due cariche nella loro posizione finale. Per portare la 3$^{\rm a}$ carica si fa il lavoro $k\big(q_1q_3/r_{13}+q_2q_3/r_{23}\big)$ contro il campo generato dalle prime due, e così via: la carica $i$--esima viene portata nella sua posizione finale nel campo delle prime $i-1$ cariche, ed il lavoro totale, e quindi l'energia della configurazione di arrivo, è
\begin{equation}
U=k\sum_{\text{coppie}}\frac{q_iq_j}{r_{ij}}=\frac12\,k\sum_{i\neq j}\frac{q_iq_j}{r_{ij}}
\tag{1.88}
\end{equation}
Il fattore 1/2 nella (1.88) è richiesto dal fatto che nella seconda sommatoria ogni coppia viene contata due volte: $i,j$ e $j,i$. Esso lo troveremo in tutte le espressioni riguardanti l'energia: anche se ha un'origine molto banale, il suo significato -- come avremo occasione di renderci conto -- non lo è.\\
La (1.88) può anche essere scritta come
\begin{equation}
U=\frac12\sum_{i=1}^nq_i\varphi^{(i)}(\vec r_i)
\tag{1.89}
\end{equation}
dove $\varphi^{(i)}(\vec r_i)$ è il potenziale nella posizione $\vec r_i$ occupata dalla carica $q_i$, dovuto a tutte le cariche salvo la $i$--esima. A questo punto ci potrebbe venir voglia di chiederci cosa ci stia a fare il fattore 1/2: non è forse vero che $q_i\varphi^{(i)}(\vec r_i)$ è l'energia potenziale dell'$i$--esima carica nel campo delle altre $n-1$, e che l'energia totale è la somma delle energie delle singole cariche? Se la risposta non viene subito alla mente, allora vuol proprio dire che quel fattore 1/2 non è del tutto banale. È vero che $q_i\varphi^{(i)}(\vec r_i)$ è l'energia dell'$i$--esima carica nel campo delle altre cariche, cioè è il lavoro che occorre fare per portare $q_i$ nel punto $\vec r_i$ quando \underline{tutte} le altre cariche sono già al loro posto, ma questo succede solo per l'ultima, ed il fattore 1/2 tiene proprio conto del fatto che il campo elettrico si instaura man mano che le cariche vengono portate nelle loro posizioni finali.\\
In generale, il fattore 1/2 non c'è quando una o più cariche vengono portate (o spostate) in un campo già costituito (\emph{campo esterno}), mentre è presente se il lavoro per spostare le cariche viene fatto contro le forze interne al sistema.\\
Consideriamo ora il caso di una distribuzione di volume che, al solito, immaginiamo costituita da tanti volumetti ``infinitesimi'' (cioè che alla fine faremo tendere a zero) $\Delta V_i$, ciascuno contenente la carica $\rho(\vec r_i)\Delta V_i$. Per la (1.89) si ha
\begin{equation}
U=\lim_{\Delta V_i\to0}\frac12\sum_i\rho(\vec r_i)\varphi(\vec r_i)\Delta V_i=\frac12\int_V\rho(\vec r\,)\varphi(\vec r\,)\,dV.
\tag{1.90}
\end{equation}
Come nella (1.89), così nella (1.90) avremmo dovuto scrivere $\varphi^{(i)}(\vec r_i)$, cioè il potenziale dovuto soltanto alle cariche esterne al volumetto $\Delta V_i$: questa volta la distinzione fra potenziale dovuto a tutte le cariche e quello dovuto solo a quelle esterne non è necessaria, infatti -- come possiamo renderci conto -- la differenza tende a zero quando $\Delta V_i\to0$: prendiamo per $\Delta V$ una sferetta di raggio $\epsilon$ sufficientemente piccolo da poter considerare costante la densità di carica al suo interno; dalla (1.70) vediamo che -- fattori numerici a parte -- il potenziale della sferetta al suo interno è dell'ordine di $\rho\Delta V/\epsilon\approx\rho\epsilon^3/\epsilon=\rho\epsilon^2$ e quindi tende a zero quando $\epsilon\to0$. Si noti che lo stesso discorso funziona per una distribuzione superficiale $\big(\sigma\epsilon^2/\epsilon\to0\big)$, ma non per una distribuzione lineare. Quindi anche per una distribuzione superficiale
\begin{equation}
U=\frac12\int_S\sigma(\vec r\,)\varphi(\vec r\,)\,dS
\tag{1.91}
\end{equation}
mentre per una distribuzione lineare l'espressione analoga alla (1.90) o (1.91) diverge: una distribuzione lineare (che però è un caso limite) non può essere costruita a partire da una distribuzione di volume e facendo tendere a zero due dimensioni, ma può essere costruita solo incollando insieme pezzettini di distribuzioni lineari, proprio come succede con le cariche puntiformi: le cariche puntiformi non le costruiamo, ma supponiamo di trovarle già belle e fatte.\\
La (1.91) è particolarmente utile nel caso dei conduttori: come vedremo, sui conduttori le cariche si dispongono solo sulla superficie, e i conduttori sono regioni equipotenziali; quindi, se $Q_i$ sono le cariche dei vari conduttori e $\varphi_i$ i rispettivi potenziali, la (1.91) diventa
\begin{equation}
U=\frac12\sum_iQ_i\varphi_i.
\tag{1.92}
\end{equation}
Nel caso delle distribuzioni di volume e di superficie, la (1.90) e la (1.91) possono (quasi sempre, cioè in tutti i casi fisicamente rilevanti) essere riscritte in un modo diverso, ma molto efficace: consideriamo il caso di una distribuzione di volume e partiamo dall'identità (che si dimostra per verifica diretta)
\begin{equation}
\dive(\varphi\vec E)=\varphi\dive\vec E+\grad\varphi\cdot\vec E=4\pi k\rho\varphi-\vec E^2
\tag{1.93}
\end{equation}
integriamo il primo e l'ultimo membro sul volume di una grande sfera di raggio $R$: il primo termine, per il teorema della divergenza dà il flusso di $\varphi\vec E$ attraverso la superficie della sfera; siccome la superficie è $O(R^2)$, se $\varphi\vec E=o(1/R^2)$ -- certamente vero se la distribuzione $\rho$ è contenuta in un volume finito $\big(\varphi=O(1/R),\ \vec E=O(1/R^2)\big)$ -- prendendo il limite per $R\to\infty$ tale flusso si annulla, per cui si ottiene
\[
4\pi k\int\rho(\vec r\,)\varphi(\vec r\,)\,dV=\int\vec E^2(\vec r\,)\,dV
\]
e quindi
\begin{equation}
U=\frac{1}{8\pi k}\int\vec E^2(\vec r\,)\,dV
\tag{1.94}
\end{equation}
dove l'integrale, anche se la distribuzione è contenuta in un volume finito, è esteso a tutto lo spazio. La (1.94) vale anche per distribuzioni superficiali; infatti una distribuzione superficiale può essere considerata come il limite di una distribuzione di volume quando il suo spessore tende a zero: nel limite il campo elettrico resta finito e quindi ha limite l'integrale a 2$^\circ$ membro della (1.94).\\
A titolo di esercizio proponiamo anche una dimostrazione diversa: supponiamo che le cariche siano distribuite su una superficie $S$; consideriamo il volume compreso fra la grande sfera di raggio $R$ e una superficie chiusa che ``fascia'' $S$, a distanza infinitesima da essa. In questo volume non sono presenti sorgenti, quindi per la (1.93) \ $\dive\varphi\vec E=-\vec E^2$. Il flusso di $\varphi\vec E$ uscente da esso, nel limite $R\to\infty$, è dato da
\fig{fig1_19}{1.19}
\[
-\int_S\varphi(\vec r\,)\Delta E_n(\vec r\,)\,dS
\]
dove $\Delta E_n$ è la differenza fra i valori della componente del campo normale ad $S$ dalle due parti della superficie, cioè è la discontinuità di $E_n$ attraverso la superficie, che sappiamo vale $4\pi k\sigma$; a questo punto la (1.94) segue immediatamente.\\
L'integrando nella (1.94), cioè $\vec E^2/8\pi k$, è detto \textbf{densità di energia}: per ora si tratta solo di una definizione; acquisterà particolare significato quando, nel caso dei campi variabili nel tempo, questa densità (assieme a quella magnetica) -- come quella di un fluido in moto non stazionario -- potrà spostarsi trasportando energia, con la velocità delle onde elettromagnetiche, da una regione ad un'altra dello spazio: per esempio dall'antenna di una stazione trasmittente alle antenne dei vari apparecchi ricevitori.

\paragraph{Esempi}
\begin{enumerate}
\item Energia di una distribuzione uniforme $\sigma$ sulla superficie di una sfera di raggio $R$.\\
Se utilizziamo la (1.91), oppure la (1.92), e la (1.69) si ha immediatamente
\begin{equation}
U=\frac12\cdot4\pi R^2\sigma\cdot k\,\frac QR=k\,\frac{Q^2}{2R}\qquad\quad Q\equiv4\pi R^2\sigma.
\tag{1.95}
\end{equation}
Un altro modo è quello di utilizzare la (1.94) ($\vec E=0$ per $r<R$)
\[
U=\frac{1}{8\pi k}\int_{r>R}\Big(k\,\frac{Q}{r^2}\Big)^2r^2\,dr\,d\Omega=\frac{4\pi}{8\pi k}\int_R^\infty k^2\,\frac{Q^2}{r^2}\,dr=k\,\frac{Q^2}{2R}\,\cdot
\]
Ma è più istruttivo ritrovare lo stesso risultato a partire dalla definizione di energia di una distribuzione di carica: supponiamo inizialmente le cariche distribuite sulla superficie di una sfera di raggio infinito, e calcoliamo il lavoro che occorre fare per contrarre la superficie fino al raggio voluto $R$. Quando il raggio della superficie è $r$, la densità superficiale è $\sigma(r)=Q/4\pi r^2$; per calcolare il lavoro per contrarre la sfera da $r$ a $r-dr$, dobbiamo determinare la componente radiale della forza che il campo esercita su ogni elemento della distribuzione superficiale. Occupiamoci preliminarmente di questo problema.\\
Consideriamo un elemento (infinitesimo) $\Delta S$ della superficie: la sua carica è $\sigma(r)\Delta S$; quant'è la forza (componente radiale) che le altre cariche esercitano sull'elemento di superficie $\Delta S$? È forse $\sigma(r)\Delta S$ moltiplicato per il campo elettrico
prodotto dalla distribuzione? E se sì, quale campo elettrico? quello sulla parte interna di $\Delta S$ (che è zero!) o quello sulla parte esterna, che vale $4\pi k\sigma(r)$? Nessuno dei due, perché quello che ci serve è il campo elettrico sulla superficie $\Delta S$ prodotto da tutte le altre cariche della distribuzione, cioè la distribuzione sferica privata dell'elemento $\Delta S$. Il campo (più precisamente la sua componente radiale) della ``sfera bucata'' nel centro del buco si determina con un bel ragionamento che fa uso del principio di sovrapposizione: la componente radiale del campo prodotto da $\Delta S$, nel limite in cui $\Delta S$ può essere considerata piana, ``per ovvie ragioni di simmetria'' (riflessione rispetto a $\Delta S$) è uguale ed opposta dalle due parti di $\Delta S$, questa si deve sommare alla componente radiale del campo dovuto a $S\setminus\Delta S$ che è continua attraverso il buco (e quindi uguale sulle due facce di $\Delta S$) per dare il campo dell'intera distribuzione, che è nullo dentro la sfera e vale $4\pi k\sigma$ subito fuori, quindi il campo cercato vale $2\pi k\sigma$. Forse non è un caso che sia la semisomma del campo dentro e fuori della sfera, prodotto dall'intera distribuzione: in effetti se avessimo considerato la distribuzione superficiale $\sigma$ come una distribuzione uniforme di volume di spessore $\Delta s$, avremmo trovato che la forza che si esercita sul tassello $\Delta S\times\Delta s$ è la carica del tassello per il valore del campo elettrico nel suo centro, che è la semisomma del campo dentro e del campo fuori.\\
Allora il lavoro che occorre fare (contro le forze del campo) per contrarre la superficie sferica dal raggio $r$ al raggio $r-dr$ è positivo (la forza esercitata dal campo è diretta verso l'esterno, indipendentemente dal segno di $\sigma$) e vale
\[
dr\int_S2\pi k\sigma(r)\cdot\sigma(r)\,dS=2\pi k\,dr\int_S\Big(\frac{Q}{4\pi r^2}\Big)^2r^2\,d\Omega=k\,\frac{Q^2}{2r^2}\,dr
\]
e quindi il lavoro per contrarre la distribuzione dall'infinito al raggio $R$ è
\[
U=k\int_R^\infty\frac{Q^2}{2r^2}\,dr=k\,\frac{Q^2}{2R}\,\cdot
\]
\item Energia di una sfera di raggio $R$ con carica Q uniformemente distribuita al suo interno.\\
Si può utilizzare la (1.90) con il potenziale dato dalla (1.70), oppure la (1.94) con il campo elettrico dato dalla (1.49), ma anche in questo caso è più istruttivo ragionare in modo diretto. Costruiamo la nostra sfera portando dall'infinito uno strato di spessore (infinitesimo) $\Delta r$ alla volta fino a formare l'intera sfera di raggio $R$. Supponiamo di avere già costruito la sfera di raggio $r$ e calcoliamo il lavoro per aggiungerci il prossimo strato, il quale contiene la carica
\begin{equation}
\Delta Q=\rho\cdot\Delta V=\frac{Q}{\frac43\pi R^3}\cdot4\pi r^2\Delta r=3Q\,\frac{r^2\Delta r}{R^3}\,\cdot
\tag{1.96}
\end{equation}
Il lavoro ha due contributi: uno quello per costruire lo strato, e questo nel limite $\Delta r\to0$ è dato dalla (1.95), e per la (1.96) è proporzionale a $(\Delta r)^2$; l'altro è il lavoro per contrarre detto strato da raggio infinito al raggio $r$ nel campo della sfera carica di raggio $r$. Questo secondo contributo è l'energia potenziale (senza il fattore 1/2!) dello strato $r$, $r+dr$ nel campo della sfera di raggio $r$, che ha carica $Qr^3/R^3$, cioè
\begin{equation}
k\,Q\,\frac{r^3}{R^3}\,\frac1r\cdot3Q\,\frac{r^2\Delta r}{R^3}=3k\,Q^2\,\frac{r^4}{R^6}\,\Delta r.
\tag{1.97}
\end{equation}
Come si vede il primo contributo (in quanto proporzionale a $(\Delta r)^2$\,) è trascurabile rispetto a quello dato dalla (1.97), che è proporzionale a $\Delta r$. Quindi sommiamo i contributi (1.97) di tutti gli strati che compongono la sfera:
\begin{equation}
U=\int_0^R3k\,Q^2\,\frac{r^4}{R^6}\,dr=k\,\frac35\,\frac{Q^2}{R}\,\cdot
\tag{1.98}
\end{equation}
Se $R\to0$ \ ($Q$ fisso), \ $U\to\infty$: una carica puntiforme richiederebbe un'energia infinita per essere costruita.
\end{enumerate}

\subsection{Le equazioni di Poisson e di Laplace}

L'equazione di Poisson (1.85) è l'equazione fondamentale dell'elettrostatica, tuttavia il nostro problema non è quello di risolverla, in quanto nell'elettrostatica delle cariche fisse la soluzione la conosciamo già, ed è data dalla (1.59). Piuttosto dobbiamo chiederci, e questo è importante anche per il seguito, se la soluzione è unica, perché se così non fosse non avremmo il diritto di considerarla come equazione fondamentale. Se $\varphi_1(x,y,z)$ e $\varphi_2(x,y,z)$ sono due soluzioni della (1.85), siccome questa è un'equazione lineare (alle derivate parziali), la loro differenza $\varphi(x,y,z)\equiv\varphi_2(x,y,z)-\varphi_1(x,y,z)$ deve essere soluzione dell'equazione omogenea associata
\begin{equation}
\Delta\,\varphi(x,y,z)=0
\tag{1.99}
\end{equation}
che è l'\textbf{equazione di Laplace}. La (1.99) certamente non ha solo la soluzione identicamente nulla: per esempio una funzione costante, oppure una funzione lineare e anche particolari polinomi di ogni grado (p.es.\ $3x_ix_j-r^2\delta_{ij}\quad\forall i,j=1\cdots3$) sono soluzioni dell'equazione di Laplace definite in tutto lo spazio (anche $\varphi=1/r$ è soluzione, ma è singolare per $r=0$). Ma l'unicità della soluzione deve essere garantita da opportune \emph{condizioni al contorno}, che sono l'analogo delle condizioni iniziali per le equazioni differenziali alle derivate ordinarie: la condizione al contorno che dobbiamo imporre alle soluzioni della (1.85), e quindi della (1.99), è -- nel caso di un sistema con carica totale finita \big(v.\ (1.78)\big) -- l'annullamento all'infinito, che corrisponde all'ipotesi fisica che non ci siano sorgenti all'infinito. Quindi dobbiamo dimostrare che l'unica soluzione dell'equazione di Laplace ovunque regolare e che si annulla all'infinito è la soluzione nulla.\\
Per dimostrare l'asserto utilizzeremo una importante proprietà delle funzioni che soddisfano (in tutto lo spazio o anche solo in una regione di esso) l'equazione di Laplace: queste funzioni sono dette \textbf{funzioni armoniche} e la regione di spazio in cui soddisfano la (1.99) è detta \emph{dominio di armonicità} della funzione. Le funzioni armoniche godono della proprietà espressa dal

\textbf{\textit{Teorema della media.}} \textit{La media di una funzione armonica sulla superficie di una sfera interamente contenuta nel dominio di armonicità è uguale al valore della funzione nel centro della sfera.}

Nell'enunciare il teorema abbiamo in mente funzioni che sono definite anche fuori del loro dominio di armonicità, come ad esempio le funzioni definite dalla (1.59), che sono armoniche solo nelle regioni prive di sorgenti: per questo abbiamo specificato che la
sfera deve essere contenuta nel dominio di armonicità.\\
Nel linguaggio dell'elettrostatica questo teorema afferma che la media del potenziale sulla superficie di una sfera che non contiene sorgenti è uguale al valore del potenziale nel centro della sfera. In altri termini, siccome la media del potenziale sulla superficie $S$ di una sfera coincide con l'energia potenziale, nel campo esterno, di una carica unitaria uniformemente distribuita sulla superficie della sfera ($\sigma=1/S$), il teorema afferma che detta energia coincide con l'energia potenziale di una carica unitaria posta nel centro.\\
La dimostrazione del teorema della media nell'ambito dell'elettrostatica, cioè nelle regioni di armonicità delle funzioni date dalla (1.59), è molto semplice ed elegante, e quindi la proponiamo a titolo di esercizio; tuttavia, siccome da un lato il teorema vale per qualsiasi funzione armonica, ma soprattutto visto che lo utilizzeremo per dimostrare l'unicità della soluzione (1.59) dell'equazione di Poisson (con le citate condizioni al contorno), non possiamo accontentarci di una dimostrazione che assume in partenza la (1.59).\\
Dimostriamo il teorema nel caso in cui $\varphi$ è il potenziale di una carica puntiforme $q$: poi il principio di sovrapposizione ci permette di estenderlo ad un numero arbitrario di cariche puntiformi, ed infine ad una distribuzione arbitraria di cariche.\\
La media di $\varphi$ sulla superficie $S$ di una sfera coincide con il lavoro per portare nella posizione finale la sfera con distribuzione superficiale di carica $\sigma=1/S$ nel campo della carica puntiforme, ma questa è anche l'energia della carica puntiforme $q$ nel campo della distribuzione superficiale, il cui potenziale all'esterno della sfera (per il teorema di Gauss) è quello di una carica unitaria nel centro della sfera:
\begin{equation}
\overline\varphi\equiv\frac1S\int_S\varphi(x,y,z)\,dS=q\cdot\frac1d
\tag{1.100}
\end{equation}
dove $d$ è la distanza fra il centro della sfera e la carica $q$, e quindi $q/d$ è il valore del potenziale $\varphi$ nel centro della sfera. $\blacksquare$

Per dimostrare il teorema in tutta la sua generalità utilizziamo la \emph{formula di Green} che si ottiene dalla (1.93): siano $f$ e $g$ due funzioni qualsiasi (naturalmente almeno di classe $C^2$); grazie all'identità espressa dalla 1$^{\rm a}$ uguaglianza della (1.93) si ha
\[
\begin{aligned}
\dive(f\grad g)&=f\Delta\,g+\grad f\cdot\grad g\\
\dive(g\grad f)&=g\Delta\,f+\grad f\cdot\grad g
\end{aligned}
\]
e sottraendo membro a membro si ottiene la formula di Green:
\begin{equation}
\dive(f\grad g-g\grad f)=f\Delta\,g-g\Delta\,f
\tag{1.101}
\end{equation}
Sia ora $\varphi$ una funzione armonica; prendiamo una sfera di raggio $R$ nel dominio di armonicità di $\varphi$ e consideriamo la funzione $1/r$, dove $r$ è la distanza dal centro della sfera: $1/r$, essendo il potenziale di una carica puntiforme, è armonica nella corona sferica di raggi $\epsilon$ e $R$. Se nella (1.101) poniamo $f=1/r$ e $g=\varphi$, nella suddetta corona sferica si ha
\begin{equation}
\dive\Big(\frac1r\grad\varphi-\varphi\grad\frac1r\Big)=\frac1r\Delta\,\varphi-\varphi\Delta\,\frac1r=0\,;\qquad\epsilon\le r\le R
\tag{1.102}
\end{equation}
in quanto $\Delta\,\varphi=0$ (per ipotesi), e $\Delta\,(1/r)=0$ per $r>0$. Quindi è nullo il flusso del
campo vettoriale $\dfrac1r\grad\varphi-\varphi\grad\dfrac1r$ \ uscente dalla corona sferica:
\begin{equation}
\begin{split}
&\frac1R\Phi_R(\grad\varphi)-\frac1\epsilon\Phi_\epsilon(\grad\varphi)+\\
&+\int_{S_R}\varphi(R,\theta,\phi)\,\frac{1}{R^2}\,dS_R-\int_{S_\epsilon}\varphi(\epsilon,\theta,\phi)\,\frac{1}{\epsilon^2}\,dS_\epsilon=0
\end{split}
\tag{1.103}
\end{equation}
I primi due termini della (1.103) sono entrambi nulli in quanto è nullo \big($\dive(\grad\varphi)=0$\big) sia il flusso di $\grad\varphi$ uscente dalla sfera di raggio $R$, sia quello uscente dalla sfera di raggio $\epsilon$; gli ultimi due integrali sono rispettivamente $4\pi$ volte la media di $\varphi$ sulla superficie di raggio $R$ e su quella di raggio $\epsilon$, che quindi per la (1.103) sono uguali. Prendendo il limite $\epsilon\to0$, si ottiene la tesi. $\blacksquare$

Una prima immediata conseguenza è che il teorema vale anche se la media è presa sul volume della sfera, e una seconda conseguenza molto importante è espressa dal seguente

\textbf{\textit{Corollario.}} \textit{Una funzione armonica non ha massimi o minimi isolati nei punti interni alla regione di armonicità.}

Infatti, se esistesse un punto di massimo (interno), esisterebbe una sfera attorno a detto punto su cui la funzione è minore che nel centro, contro il teorema della media. Quindi una funzione armonica ha i suoi massimi e i suoi minimi solo sulla frontiera del dominio di armonicità.\\
Tradotto nel linguaggio dell'elettrostatica ciò significa che in un campo elettrico, in una regione priva di sorgenti non esistono posizioni di equilibrio stabile per le cariche. Certamente possono esistere punti in cui il campo elettrico si annulla, ma in un intorno di detti punti si hanno sia linee di forza entranti che uscenti (v.\ figura 1.1), come è richiesto dal fatto che il flusso del campo elettrico uscente da un qualsiasi volume privo di sorgenti deve essere nullo: in effetti la non esistenza di massimi o minimi può essere vista come conseguenza del teorema della divergenza. Si noti anche che il risultato non è più vero nelle regioni in cui sono presenti le sorgenti: per esempio il centro di una sfera uniformemente carica è un punto di equilibrio stabile per le cariche di segno opposto a quello della distribuzione.

Torniamo al problema dell'unicità della soluzione della (1.99) che si annulla all'infinito: siccome può avere i suoi massimi e i suoi minimi solo all'infinito dove è nulla, essa deve essere nulla ovunque. Il risultato può essere generalizzato al caso in cui la regione di armonicità non è tutto lo spazio e la sua frontiera è costituita da diverse superfici (per esempio le superfici di un certo numero di conduttori): è unica la funzione armonica che ha assegnati valori su tutta la frontiera della regione in cui è definita. Infatti se due funzioni armoniche obbediscono alle stesse condizioni al contorno, la loro differenza è nulla su tutta la frontiera e quindi, come sopra, è nulla dappertutto.

\bigskip
\noindent\rule{5cm}{0.4pt}

\textbf{Nota.} A partire dal prossimo capitolo tutte le equazioni saranno scritte nel sistema di Gauss, quindi non comparirà più il fattore $k$ in quanto $k=1$.
